% concatenated from arxiv.tex
 \documentclass[a4paper,12pt,reqno]{amsart}
%\usepackage{times} %,fullpage}%,backref}
\usepackage{mathrsfs,latexsym, bm, amscd,amssymb,url}
\pagestyle{headings}
\usepackage{xcolor, mathtools}
\PassOptionsToPackage{pdfusetitle,pagebackref,colorlinks}{hyperref}
\usepackage[all,cmtip]{xy}
\usepackage{bookmark}
\hypersetup{
linkcolor={red!70!black},
citecolor={green!70!black},
urlcolor={teal!80!black},
linkcolor={teal!80!black}
}
%\definecolor{dblue}{rgb}{0,0,.6}
%\usepackage[colorlinks=true, linkcolor=dblue, citecolor=dblue, filecolor = dblue, menucolor = dblue, urlcolor = dblue]{hyperref}

\textwidth=450pt 
\oddsidemargin=12pt
\evensidemargin=12pt
\renewcommand{\baselinestretch}{1.2}
\setlength{\footskip}{25pt}
%
%\setlength{\headheight}{6.15pt}
%\setlength{\headsep}{0.5cm}
%\setlength{\headheight}{6.15pt}
%\setlength{\headsep}{1.0cm}
%\usepackage[latin1]{inputenc} 	
%\usepackage[T1]{fontenc}         
%\usepackage{ae}									
%\usepackage[arrow, matrix, curve]{xy}  
\usepackage[all,cmtip]{xy}  
\usepackage{graphicx}
%\usepackage{MnSymbol} 
\usepackage{color}
\usepackage{hyperref}
\usepackage{enumerate}
\usepackage{stmaryrd}



\usepackage[T1]{fontenc}         
\usepackage{ae}									
%\usepackage[arrow, matrix, curve]{xy}  
\usepackage[all,cmtip]{xy}  
\usepackage{graphicx}
%\usepackage{MnSymbol} 
\usepackage{color}

\interfootnotelinepenalty=10000


%----------------------------------------------
% Theorem like environments
%


\newtheorem{theorem}{Theorem}[section]
\newtheorem*{main}{Main Theorem}
\newtheorem{lemma}[theorem]{Lemma}
\newtheorem{proposition}[theorem]{Proposition}
\newtheorem{corollary}[theorem]{Corollary}
\newtheorem{claim}[theorem]{Claim}
\newtheorem{conjecture}[theorem]{Conjecture}
%\newtheorem{question}[theorem]{Question}
\newtheorem{example}[theorem]{Example}
\newtheorem{step}{Step}[subsection]
\renewcommand{\thestep}{\arabic{step}}

\newtheorem{alphtheorem}{Theorem}
\renewcommand{\thealphtheorem}{\Alph{alphtheorem}}
\newtheorem{alphtheoremprime}{Theorem}
\renewcommand{\thealphtheoremprime}{\Alph{alphtheoremprime}'}







\theoremstyle{question}
\newtheorem{question}[theorem]{Question}
\theoremstyle{definition}
\newtheorem{definition}[theorem]{Definition}
\newtheorem{notation}[theorem]{Notation}

\theoremstyle{remark}
\newtheorem{remark}[theorem]{Remark}

\newtheoremstyle{cited}{.5\baselineskip\@plus.2\baselineskip\@minus.2\baselineskip}{.5\baselineskip\@plus.2\baselineskip\@minus.2\baselineskip}{\itshape}{}{\bfseries}{\bfseries .}{5pt plus 1pt minus 1pt}{\thmname{#1}\thmnumber{~#2}\thmnote{ \normalfont#3}}
\theoremstyle{cited}
\newtheorem{citedthm}[theorem]{Theorem}
\newtheorem{citedconj}[theorem]{Conjecture}
\newtheorem{citedlem}[theorem]{Lemma}
\newtheorem{citedprop}[theorem]{Proposition}

\newtheoremstyle{citeddef}{.5\baselineskip\@plus.2\baselineskip\@minus.2\baselineskip}{.5\baselineskip\@plus.2\baselineskip\@minus.2\baselineskip}{}{}{\bfseries}{\bfseries .}{5pt plus 1pt minus 1pt}{\thmname{#1}\thmnumber{~#2}\thmnote{ \normalfont#3}}
\theoremstyle{citeddef}
\newtheorem{citednot}[theorem]{Notation}


\setcounter{tocdepth}{1}

%---------------------------------------------------


%-----------
%-----------Yagna's typsets--------
\usepackage{tikz-cd}
\usepackage{enumitem}

\PassOptionsToPackage{pdfusetitle,pagebackref,colorlinks}{hyperref}
\usepackage{bookmark}
\hypersetup{
linkcolor={red!70!black},
citecolor={green!70!black},
urlcolor={blue!80!black}
}


\def\Z{{\bf Z}}
\def\N{{\bf N}}
\def\F{{\bf F}}
\def\C{{\bf C}}
\def\A{{\bf A}}
\def\R{{\bf R}}
\def\Q{{\bf Q}}
\def\P{{\bf P}}
\def\T{{\bf T}}
\def\D{{\bf D}}
\def\k{{\bf k}}


\def\CC{{\mathbb C}}
\def\QQ{{\mathbb Q}}
\def\PP{{\mathbb P}}
\def\ZZ{{\mathbb Z}}
\def\HH{{\mathbb H}}


\def\AT{(A,\theta)}
\def\CT{{\mathcal CT}_5}
\def\JT{(JC,\theta)}
\def\XT{(X,\Theta)}

\def\av{abelian variety}
\def\avs{abelian varieties}
\def\ppav{principally polarized abelian variety}
\def\pps{principally polarized abelian surface}
\def\ppavs{principally polarized abelian varieties}
\def\pav{polarized abelian variety}
\def\pavs{polarized abelian varieties}

\def\ndeg{nondegenerate}
\def\gndeg{geometrically nondegenerate}

\def\lra{\longrightarrow}
\def\lraa{\lra\hskip-6.05mm \lra}
\def\llra{\hbox to 10mm{\rightarrowfill}}
\def\llraa{\llra\hskip-9.2mm \llra}
\def\lllra{\hbox to 15mm{\rightarrowfill}}
\def\lllraa{\lllra\hskip-14.2mm \lllra}

\def\mapright#1{\smash{\mathop{\longrightarrow}\limits^{#1}}}
\def\mapstoright#1{\smash{\mathop{\longmapsto}\limits^{#1}}}
\def\mapdown#1{\big\downarrow
\rlap{$\vcenter{\hbox{$\scriptstyle#1$}}$}}
\def\mapup#1{\big\uparrow
\rlap{$\vcenter{\hbox{$\scriptstyle#1$}}$}}
\def\mapdownl#1{\llap{$\vcenter{\hbox{$\scriptstyle#1$}}$}\big\downarrow}
\def\mapupl#1{\llap{$\vcenter{\hbox{$\scriptstyle#1$}}$}\big\uparrow}
\def\ww{\bigwedge\nolimits}
\def\pp{\partial}

\def\PA{{\widehat A}}
\def\PB{{\widehat B}}
\def\PK{{\widehat K}}
\def\PV{{\widehat V}}
\def\PX{{\widehat X}}
\def\PE{{\widehat E}}
\def\PF{{\widehat F}}
\def\PY{{\widehat Y}}
\def\PS{{\widehat S}}
\def\PT{{\widehat T}}
\def\PC{{\widehat C}}
\def\Pf{{\widehat f}}
\def\Pg{{\widehat g}}
\def\Pp{{\widehat\pi}}
\def\pp{{\widehat p}}

\def\phi{{\varphi}}
\def\wf{{\widetilde f}}
\def\wX{{\widetilde X}}
\def\wL{{\widetilde L}}
\def\wZ{{\widetilde Z}}
\def\wY{{\widetilde Y}}
\def\wA{{\widetilde A}}


\def\bP{\mathbf{P}}


\def\cI{\mathcal{I}}
\def\cD{\mathcal{D}}
\def\cA{\mathcal{A}}
\def\cF{\mathcal{F}}
\def\cL{\mathcal{L}}
\def\cO{\mathcal{O}}
\def\cG{\mathcal{G}}
\def\cP{\mathcal{P}}
\def\cH{\mathcal{H}}
\def\cEx{\mathcal{Ext}}

\def\cE{\mathcal{E}}
\def\cS{\mathcal{S}}
\def\cT{\mathcal{T}}
\def\cC{\mathcal{C}}
\def\cM{\mathcal{M}}
\def\cN{\mathcal{N}}
\def\cR{\mathcal{R}}
\def\cJ{\mathcal{J}}
\def\cX{\mathcal{X}}
\def\cQ{\mathcal{Q}}
\def\cB{\mathcal{B}}
\def\cZ{\mathcal{Z}}
\def\cV{\mathcal{V}}
\def\cW{\mathcal{W}}


\let\tilde\widetilde

\def\dra{\dashrightarrow}
\def\tto{\twoheadrightarrow}
\def\isom{\simeq}
\def\eps{\varepsilon}
\def\ot{\otimes}
\def\ie{\hbox{i.e.,}}
\def\eg{\hbox{e.g.,}}
\def\irr{irreducible}



\def\Id{\hbox{\rm Id}}
\def\Im{\mathop{\rm Im}\nolimits}
\def\vide{\varnothing}
\def\k{\mathbf k}



\DeclareMathOperator{\Alb}{Alb}
\DeclareMathOperator{\alb}{alb}
\DeclareMathOperator{\Aut}{Aut}
\DeclareMathOperator{\Bl}{Bl}
\DeclareMathOperator{\Bs}{Bs}
\DeclareMathOperator{\ch}{ch}
\DeclareMathOperator{\CH}{CH}
\DeclareMathOperator{\Coh}{Coh}
\DeclareMathOperator{\codim}{codim}
\DeclareMathOperator{\cSpec}{\mathbf{Spec}}
\DeclareMathOperator{\dD}{D}
\DeclareMathOperator{\ddeg}{deg}
\DeclareMathOperator{\diff}{Diff}
\DeclareMathOperator{\Ext}{Ext}
\DeclareMathOperator{\Exc}{Exc}
\DeclareMathOperator{\Hom}{Hom}
\DeclareMathOperator{\id}{id}
\DeclareMathOperator{\Ker}{Ker}
\DeclareMathOperator{\lc}{lc}
\DeclareMathOperator{\PGL}{PGL}
\DeclareMathOperator{\Pic}{Pic}
\DeclareMathOperator{\PProj}{\mathbf{Proj}}
\DeclareMathOperator{\Ram}{Ram}
\DeclareMathOperator{\rank}{rank}
\DeclareMathOperator{\Sing}{Sing}
\DeclareMathOperator{\Spec}{Spec}
\DeclareMathOperator{\Supp}{Supp}
\DeclareMathOperator{\Sym}{Sym}
\DeclareMathOperator{\td}{td}
\DeclareMathOperator{\Td}{Td}
\DeclareMathOperator{\tr}{Tr}
\DeclareMathOperator{\vol}{vol}
%----------------


\begin{document}

\title[Geometry of Holomorphic One-forms on Smooth Projective Varieties]{Geometry of Holomorphic One-forms on Smooth Projective Varieties}
\author{Jiabin Du}
\address{Shanghai Institute for Mathematics and Interdisciplinary Sciences (SIMIS), Shanghai 200433, China}
\address{Research Institute of Intelligent Complex Systems, Fudan University, Shanghai 200433, China}
\email{jiabin.du@simis.cn}

\author{Feng Hao}
\address{School of Mathematics\\ Shandong University\\ Jinan 250100\\ China}
\email{feng.hao@sdu.edu.cn}

\author{Haoyuan Li}
\address{School of Mathematics\\ Shandong University\\ Jinan 250100\\ China}
\email{haoyuan.li@mail.sdu.edu.cn}

\author{Zichang Wang}
\address{School of Mathematical Science\\ University of Science and Technology of China\\ Hefei 230026\\ China}
\email{wangzichang@mail.ustc.edu.cn}



\date{\today}

\keywords{holomorphic one-forms, zero loci, abelian variety, linearity,  without blow-up in codimension zero, cohomology jump loci}

\subjclass{14J17, 14K05, 14K12}
 

\begin{abstract}
%In this article, we show  any given morphism from a smooth projective varieties $X$ to a simple abelian varieties $A$ are smooth if and only if there exists a nowhere vanishing holomorphic 1-form on $X$ that  is pulled back from $A$.  Moreover, we systematically studied the (non)linearity of the spaces of degenerate holomorphic 1-forms on smooth projective varieties.

%In this article, we show  that any  morphism from a smooth projective variety $X$ to a simple abelian variety $A$ is smooth, if and only if there exists a nowhere vanishing holomorphic 1-form on $X$ pulled back from $A$. As the key ingredient in the proof, we also show any $\mathbb{Z}$-homological fibre bundle morphism is without blow-up in codimension 0.


In this article, we show  that any  morphism $f$ from a smooth projective variety $X$ to a simple abelian variety $A$ is smooth, if and only if there exists a holomorphic 1-form $\omega$ on $A$  such that $f^*\omega$ has no zero. As the key ingredient in the proof, we show any $\mathbb{Z}$-homology fibre bundle morphism is without blow-up in codimension 0 in the sense of Sabbah.
%As a key ingredient in the proof, we also estabilish that any $\mathbb{Z}$-homological fibre bundle morphism is without blow-up in codimension 0.
%This proves a conjecture proposed by the second named author.

Furthermore, we investigate the structure of the spaces of holomorphic 1-forms with zeros, and show that they are linear for large classes of varieties. Also, we construct a delicate example of a smooth projective subvariety of an abelian variety for which the holomorphic 1-forms with positive dimensional zero loci do not form a linear subset. Finally, we study algebraic surfaces admitting holomorphic 1-forms that have zeros and do not arise from cohomology jump loci.

%In addition, we investigate the linearity of the spaces of degenerate holomorphic 1-forms, construct an explicit example of a smooth subvariety of an abelian variety with a non-linear lower level set and analyze the structure of algebraic surfaces admitting degenerate holomorphic 1-forms  not arising from cohomology jump loci. 
\end{abstract}
\maketitle
\vspace*{6pt}
\tableofcontents  % for this guide only.
% A table of contents should normally not be included
\section{Introduction}

Holomorphic 1-form is a basic concept in algebraic geometry and occupies a central position in analyzing the geometry and topology of  irregular varieties. The space of holomorphic 1-forms governs the morphisms from varieties to abelian varieties, especially the Albanese maps, so the information encoded in holomorphic 1-forms should potentially reveal many structural properties of irregular varieties. In fact, recent studies  (\cite{CCH25}, \cite{Chu24}, \cite{DHL24}, \cite{HK05}, \cite{Hao24}, \cite{HS21}, \cite{Kot22}, \cite{LZ05},  \cite{Pie25}, \cite{PS14}, \cite{Sch21}, \cite{SY25}, etc.) show that fine properties of holomorphic 1-forms heavily constrain the topology and birational geometry of irregular varieties. 

Since the zeros of holomorphic 1-forms are controlled by the singularities of the Albanese maps, it is usually very difficult to keep track of the changes of zeros of holomorphic 1-forms under algebraic operations on varieties like birational modifications, branched covers, base changes, and  hyperplane sections, etc. In this article, we provide a fine study of the space of holomorphic 1-forms without or with zeros on smooth projective varieties. Throughout this article, all varieties are defined over the complex numbers field $\mathbb{C}$. 

We develop a new method via limit relative conormal vectors to analyze the singularities of morphisms from smooth projective varieties to abelian varieties. For our first main theorem, we have
\begin{theorem}[{=Theorem \ref{thm:simple-nowhere vanishing 1-form}}]\label{thm:simple-smooth}
    Let $f\colon X\to A$ be a morphism from a smooth projective variety $X$ to a simple abelian variety $A$. The following are equivalent:

    1) $f$ is smooth;
    
    2) there exists a holomorphic 1-form $\omega\in H^0(A, \Omega^1_A)$ such that $f^*\omega$ has no zero.
\end{theorem}

The simpleness condition in Theorem \ref{thm:simple-smooth}  is necessary, and Theorem \ref{thm:simple-smooth} is sharp.  In fact, Schreieder and Yang \cite{SY25} provide the following example. Let $Y=E_1\times E_2\times \mathbb{P}^1$, with $E_1$ and $E_2$ being two non-isogenous elliptic curves. Denote by $C$ the disjoint union of two elliptic curves $E_1\times \{e_2\}\times \{0\}$ and $\{e_1\}\times E_2\times \{\infty\}$. Consider the blowup $\text{Bl}_{C}Y$ of $Y$ along $C$. Then the natural morphism $\pi\colon \text{Bl}_{C}Y \to E_1\times E_2$ is not smooth. However, there exists a nowhere vanishing holomorphic 1-form pulled back from $E_1\times E_2$. Theorem \ref{thm:simple-smooth}  completely proves a conjecture that is proposed by the second named author in \cite[Conjecture 1.5]{Hao24}. For varieties of dimension no more than 4, this problem is already solved by Church \cite[Theorem D]{Chu26}. 



In order to deal with Theorem \ref{thm:simple-smooth}, we establish our second main result. Recall the following key terminology introduced by Sabbah \cite[Definition 3.1]{Sab85}. A proper morphism $f\colon X\to Y$ between smooth varieties is \textit{without blow-up in codimension 0}, if the natural morphism from the relative conormal space of $f$ to $Y$ is equidimensional. Here we mean the relative conormal space of $f$ by the closure of the relative conormal bundle of the smooth loci of $f$ in the cotangent bundle of $X$ (See Definition \ref{sans  \'eclatement en codimension 0} in Section \ref{sect:smooth-mor-simpleAV-1-form}). 

\begin{theorem}[{=Theorem \ref{thm:homological-trivial-deformation-no-blow-up-codim0}}]\label{main-thm:homological-trivial-deformation-no-blow-up-codim0}
  Let $f\colon X\to Y$ be a proper morphism between algebraic varieties, where $X$ is smooth and $Y$ is normal. Suppose $R^if_*\mathbb{Z}$ are $\mathbb{Z}$-local systems on $Y$ for all integers $i$, then $Y$ is smooth and $f$ is without blow-up in codimension 0.
\end{theorem}

The morphism $f$ with the topological condition that $R^if_*\mathbb{Z}$ are $\mathbb{Z}$-local systems for all integers $i$, is usually called a \textit{$\mathbb{Z}$-homology fibre bundle}. Roughly speaking, Theorem \ref{main-thm:homological-trivial-deformation-no-blow-up-codim0} shows that the $\mathbb{Z}$-homology fibre bundle condition kills all the hidden non-flatness of a morphism. For the morphism $f\colon X\to A$ in Theorem \ref{thm:simple-smooth}, via the simpleness of $A$ and Theorem \ref{main-thm:homological-trivial-deformation-no-blow-up-codim0}, we  obtain that the space of holomorphic 1-forms with zeros on $X$ is either the zero vector space or $H^0(X, \Omega_X^1)$. Then the existence of a nowhere vanishing holomorphic 1-form on $X$ pulled back from $A$ induces the smoothness of $f$. Theorem \ref{main-thm:homological-trivial-deformation-no-blow-up-codim0} itself is also motivated by the de Bobadilla-Koll\'ar conjecture, which states that if $f\colon X\to Y$ is a proper morphism from a smooth variety $X$ to a normal variety $Y$, and $R^if_*\mathbb{Z}$ are $\mathbb{Z}$-local systems on $Y$ for all $i$, then $f$ is a smooth morphism (See \cite[Conjecture 3]{Bobadilla-Kollar2012}). Note that this conjecture is still wildly open. It is even not known for morphisms of relative dimension 2.

Theorem \ref{thm:simple-smooth} shows that if the Albanese variety of a variety $X$ is simple, then the space of holomorphic 1-forms with zeros on $X$ is either the zero vector space or the whole vector space of holomorphic 1-forms. For a general smooth projective variety $X$, we aim to describe the structure of the space of holomorphic 1-forms
with zeros and its lower level sets:
\[
W^i(X):=\{\omega\in H^0(X, \Omega_X^1)\ |\ \codim Z(\omega)\leq i \},
\] 
where $Z(\omega)$ is the set theoretic zero loci of $\omega$. These lower level sets, introduced by Budur-Wang-Yoon \cite{BWY16}, are all algebraic cones nailed to the origin of $H^0(X,\Omega_X^1)$. Each lower level set $W^i(X)$ contains \textit{linear subsets}\footnote{We say a subset of a complex vector space a \textit{linear subset}, if it is a union of finitely many vector subspaces.} $V^i(X)$ (see Section \ref{sect:deg-1-form-linear} for the definition) consisting of holomorphic 1-forms arising from cohomology jump loci. The linearity of $V^i(X)$ is essentially coming from the formality of the variety $X$. However, the structure of $W^i(X)$ is much more delicate and encodes much finer algebraic nature of the variety $X$. For the structure of the whole space of holomorphic 1-forms with zeros, we prove the following theorem.

\begin{theorem}[{=Theorem \ref{thm:linearity-zeros-gmm}, \ref{prop:simple-linear}, \ref{prop:23linearity}}]\label{prop:gmm-sav-linear}
    Let $X$ be a smooth projective variety of dimension $n$. If either
    
    1) $X$ is a good minimal model, or

    2) $X$ has simple Albanese variety, or

    3) $n\leq 3$,
    
    \noindent then the space $W^n(X)$ of holomorphic 1-forms with zeros is a union of finitely many vector subspaces of $H^0(X,\Omega_X^1)$. 
\end{theorem}

Theorem \ref{prop:gmm-sav-linear} with assumption 1) can be deduced from the main result of \cite{HWZ26}, and Theorem \ref{prop:gmm-sav-linear} with assumption 2) is a consequence following from Theorem \ref{thm:simple-smooth}. Theorem \ref{prop:gmm-sav-linear} with assumption 3) follows from the classification of surfaces and threefolds with nowhere vanishing holomorphic 1-forms \cite{Sch21} and \cite{HS21}. In general, one could expect that at least the birational modification would make the structure of the space of holomorphic 1-forms with zeros rather delicate and out of control. The structure of $W^n(X)$ relies on a concrete description of the birational operations involved in the minimal model program from the smooth projective variety $X$ to its minimal models.

As for the lower level sets, Wang-Zhang \cite{WZ25} provides an example of a threefold $X$ arising from a careful blowup of a three dimensional subvariety $Y$ of an abelian fourfold along a general curve in $Y$ such that $W^2(X)$ is not linear. However, for smooth projective subvarieties $X$ of abelian varieties, all the lower level sets $W^i(X)$ were expected to be linear by experts. Surprisingly, we construct an explicit example of a smooth projective fourfold $X$, which is a subvariety of an abelian variety, such that $W^3(X)$ is not linear.

\begin{theorem}[{=Theorem \ref{thm:non-linear example}}]\label{main-thm:nonlinear-ex}
There exists a four dimensional smooth projective subvariety $X$ of an abelian sixfold, such that the third lower level set $W^3(X)$ is not linear.
\end{theorem}

Theorem \ref{main-thm:nonlinear-ex} answers a question \cite[Question 1.10]{DHL24}. Analyzing the non-linearity of $W^3(X)$ is rather delicate, but the construction of $X$ in Theorem \ref{main-thm:nonlinear-ex} can be easily described as follows. Let $A$ be a simple abelian 5-fold and $E$ be an elliptic curve. Take a general hyperplane section $Y$ of $A$ and a general smooth curve $C$ in $Y$. We carefully choose a hyperplane section $X$ of $Y\times E$ that passes through $C\times E$ such that $X$ is smooth and in a nice position. Then $X$ is a smooth subvariety of $A\times E$, and one can show $W^3(X)$ is not linear. 

The example in Theorem \ref{main-thm:nonlinear-ex} shows that the holomorphic 1-forms with zeros that are not from cohomology jump loci are mysterious and should reflect deep algebro-geometric nature of varieties. In the last part of this article, we provide a cohomological criterion for this kind of holomorphic 1-forms, which is easily deduced from \cite{GL87}. Also, we study smooth projective surfaces admitting holomorphic 1-forms that have zeros and do not come from cohomology jump loci.



\noindent \textbf{Acknowledgment} We would like to thank Nero Budur, Yongqiang Liu, Stefan Schreieder, and Botong Wang for very useful conversations and comments, and also thank the Shanghai Institute for Mathematics and Interdisciplinary Sciences (SIMIS) for providing an excellent research environment for discussions related to this work. The second named author is supported by NSFC No.1240010723, SDNSFC (No. 2024HWYQ-009, No. ZR2024MA007, No. tsqn202312060).


\section{Limit relative conormal vectors and morphism to simple abelian varieties}\label{sect:smooth-mor-simpleAV-1-form}

In this section, we provide a fine study on holomorphic 1-forms and morphisms from smooth projective varieties to simple abelian varieties.

Let $f\colon X\to Y$ be a proper morphism between smooth varieties, and $X^\circ$ be the maximal open subset of $X$ such that $f|_{X^\circ}$ is smooth.
Denote the cotangent bundle of $X$ by $T^*X$ and the scheme theoretic fibre of $f$ over $y$ by $F_y$ for any closed point $y\in Y$. The \textit{relative conormal  space} $T_f^*X$ is defined to be the closure of the subvariety $\{ (x,\eta)\in T^*X|_{X^\circ} \  | \  \eta(T_xF_{f(x)} ) = 0   \}$ in $T^*X$, where $T_xF_{f(x)}$ is the tangent space of the fibre $F_{f(x)}$ at the closed point $x$. Recall the following concept introduced by Sabbah \cite[Definition 3.1]{Sab85}.


\begin{definition}\label{sans  \'eclatement en codimension 0}
We say that a proper morphism $f\colon X\to Y$ between smooth varieties is \textit{without blow-up  in codimension 0}, if the natural composition morphism $T^*_fX\to X\to Y$ is equidimensional.
\end{definition}



\begin{theorem}\label{thm:homological-trivial-deformation-no-blow-up-codim0}
  Let $f\colon X\to Y$ be a proper morphism between algebraic varieties, where $X$ is smooth and $Y$ is normal. Suppose $R^if_*\mathbb{Z}$ are $\mathbb{Z}$-local systems on $Y$ for all integers $i$, then $Y$ is smooth and $f$ is without blow-up in codimension 0.
\end{theorem}

\begin{proof}
The smoothness of $Y$ is due to de Bobadilla-Koll\'ar \cite[Corollary 7]{Bobadilla-Kollar2012}. Again by \cite[Corollary 7]{Bobadilla-Kollar2012} and \cite[Proposition 6]{Bobadilla-Kollar2012}, we have $f$ is flat, and each scheme theoretic fibre is generically reduced, irreducible and a local complete intersection. Hence each scheme theoretic fibre is integral. Then $f$ is smooth at a closed point $x\in X$ if and only if $x$ is a smooth closed point of the fibre $F_{f(x)}$ by \cite[Proposition 6]{Bobadilla-Kollar2012}, i.e., $X-X^\circ$ is exactly the set of singular points of singular fibres of $f$. 

Consider the composition map $\pi\colon T_f^*X\to X\to Y$. We need to show $\pi$ is equidimensional. Denote the discriminant of $f$ by $$\Delta(f) \coloneqq \{y\in Y\ |\  f\ \text{is singular along some point in}\ f^{-1}(y)\}.$$ The previous discussion shows $$\Delta(f) = \{y\in Y\ |\ F_y\ \text{is singular}\}.$$

For any closed point $y\in Y-\Delta(f)$, it is clear that  the fibre $\pi^{-1}(y)$ is exactly the total space of the conormal bundle $\mathcal{N}^*_{F_y/X}$, and $\dim \pi^{-1}(y)=\dim \mathcal{N}^*_{F_y/X}=\dim X$. Fix a closed point $y_0 \in \Delta(f)$, then $F_{y_0}$ is singular. Denote the nonempty smooth locus of $F_{y_0}$ by $F^\circ_{y_0}$. Then by the previous discussion, $F^\circ_{y_0}=X^\circ\cap F_{y_0}$. By the definition of relative conormal space, it is clear that $$T^*_fX|_{F^\circ_{y_0}}=\mathcal{N}^*_{F^\circ_{y_0}/X}.$$
Since $F_{y_0}$ is integral, $T^*_fX|_{F^\circ_{y_0}}$ is  irreducible, so is the closure $\overline{T_f^* X|_{F^\circ_{y_0}}}^{T^*X}$ of  $T_f^* X|_{F^\circ_{y_0}}$ in $T^*X$. Also, $$\dim \overline{T_f^* X|_{F^\circ_{y_0}}}^{T^*X}=\dim T_f^* X|_{F^\circ_{y_0}}=\dim X.$$ Hence to show $\pi$ is equidimensional, it suffices to show
$\pi^{-1}(y_0) = \overline{T_f^* X|_{F^\circ_{y_0}}}^{T^*X}$. Note that $T_f^* X|_{F^\circ_{y_0}}\subset \pi^{-1}(y_0)$ and $\pi^{-1}(y_0)$ is closed in both $T^*_fX$ and $T^*X$. Thus we just need to show   $\pi^{-1}(y_0) = \overline{T_f^* X|_{F^\circ_{y_0}}}^{\pi^{-1}(y_0)}$, i.e., to show $T_f^* X|_{F^\circ_{y_0}}$ is dense in $\pi^{-1}(y_0)$.  To this end, for any fixed closed point $x_0\in F_{y_0}-F^\circ_{y_0}$, we only need to prove for any Cauchy sequence $\{(a_k,\eta_k)\}_{k\in \mathbb{Z}^+}$ in $T_f^*X|_{X^\circ}$ with limit point $(x_0, \eta_0)\in T^*X$, there exists a Cauchy sequence $\{(b_k,\mu_k)\}_{k\in \mathbb{Z}^+}$ in $T_f^*X|_{F^\circ_{y_0}}$ which has the same limit point $(x_0, \eta_0)\in T^*X$, as is described in Figure \ref{fig:fibration}. 
\begin{figure}[htbp]
\centering
\includegraphics[width=0.8\textwidth]{figure1.png}
\caption{}\label{fig:fibration}
\end{figure}






Assume $\dim Y = m$ and $\dim X=n$, take a small open ball $U_{y_0}$ together with complex analytic coordinates $z_1, \ldots, z_m$ around $y_0$. Then $dz_1, \ldots, dz_m$ are nowhere vanishing pointwise linearly independent local holomorphic 1-forms on $U_{y_0}$. For any closed point $x\in X^\circ\cap f^{-1}(U_{y_0})$ at which $f$ is smooth, the cotangent map $f^*_x\colon T^*_{f(x)}Y\to T^*_xX$ is injective (the dual of the surjective tangent map at $x$). Thus $$f^*_x(T_{f(x)}^*Y)=\text{Span}\langle f^*dz_1(x),...,f^*dz_m(x)\rangle,$$ and $\dim f^*_x(T_{f(x)}^*Y)=m$. Moreover, we have $$f^*_x(T_{f(x)}^*Y)\subseteq\{\alpha\in T^*_xX\ |\ \alpha(T_xF_{f(x)})=0\}=\text{the fibre of\ }T_f^*X|_{X^{\circ}}\ \text{over} \ x.$$ Since $f$ is flat, $\dim T_xF_{f(x)}=n-m$. Hence $\dim \{\alpha\in T^*_xX\ |\ \alpha(T_xF_{f(x)})=0\}=m$ as $\dim T^*_xX=n$.
Therefore, we get 
\begin{equation}\label{eq:fibre-localcoord-represent}
  \text{the fibre of\ }T_f^*X|_{X^{\circ}}\ \text{over} \ x=f^*_x(T_{f(x)}^*Y)=\text{Span}\langle f^*dz_1(x),...,f^*dz_m(x)\rangle.  
\end{equation}
Given any Cauchy sequence $\{(a_k,\eta_k)\}_{k\in \mathbb{Z}^+}$ in $T_f^*X|_{X^\circ \cap f^{-1}(U_{y_0})}$ with limit point $(x_0, \eta_0)\in T^*X$. Since $a_k\in X^\circ \cap f^{-1}(U_{y_0})$, we have $\eta_k \in \text{span}\langle f^*dz_1(a_k),...,f^*dz_m(a_k)\rangle$ by Equation (\ref{eq:fibre-localcoord-represent}). Denote $\eta_k = \sum_{i=1}^m \alpha_i^k f^*dz_i(a_k)$, where $\alpha_i^k\in\mathbb{C}$. Now consider the sequence $\{(b_k,\mu_k)\}$ defined as follows: Choose $b_k$ to be any Cauchy sequence converging to $x_0\in F_{y_0}-F^\circ_{y_0}$ such that $b_k\in F^{\circ}_{y_0}$. Define $$\mu_k = \sum_{i=1}^m \alpha_i^k f^*dz_i(b_k),$$ then $$\mu_k\in f^*_{b_k}(T^*_{y_0}Y)=\text{the fibre of\ }T_f^*X|_{X^{\circ}}\ \text{over} \ b_k=\text{the fibre of\ }T_f^*X|_{F^\circ_{y_0}}\ \text{over} \ b_k$$
Thus $(b_k,\mu_k)\in T_f^*X|_{F^\circ_{y_0}}$ for any $k$. Since $\{f^*dz_1, \ldots, f^*dz_m\}$ are holomorphic 1-forms on $f^{-1}(U_{y_0})$, we conclude by continuity that the sequence $\{(b_k,\mu_k)\}_{k\in\mathbb{Z}^+}$ is indeed a Cauchy sequence and has the same limit point $(x_0, \eta_0)\in T^*X$ as $\{(a_k,\eta_k)\}_{k\in\mathbb{Z}^+}$.
\end{proof}

%\begin{theorem}\label{thm:equidimensional}
 %   Let $f\colon X\to A$ be a morphism from a smooth projective variety $X$ to a simple abelian variety $A$, $\phi_f\colon  T_f^*X \to X$ be the relative conormal space on $X$,  If there exists a holomorphic 1-form $\omega\in H^0(A, \Omega^1_A)$ such that $f^*\omega$ has no zero, then the composition map $\pi_f = f\circ \phi_f \colon   T_f^*X \to A$ is equidimensional.
%\end{theorem}

%\begin{proof}
 %   By Duta-Hao-Liu {\color{red}\cite{DHL24}}, If there exists a holomorphic 1-form $\omega\in H^0(A, \Omega^1_A)$ such that $f^*\omega$ has no zero, then $R^if_*\mathbb{Z}$ are $\mathbb{Z}$-local systems for all $i$. And Bobadilla-Koll\'ar {\color{red}\cite{Bobadilla-Kollar2012}} proved if $g\colon X\to Y$ is a morphism between smooth varities, and $R^ig_*\mathbb{Z}$ are $\mathbb{Z}$-local systems for all $i$, then 
  %  \begin{itemize}
     
   % \item[(1)] $g$ is smooth at point $x\in X$ if and only if $x$ is  smooth  on fiber $F_{g(x)}$ .
    %\item[(2)] $g$ is a flat morphism and each fiber of $g$ is integral.
   % \end{itemize}
  %  So for our setting, $X-X^\circ$ is exactly the set of singular points of singular fibers of $f$. Denote $\Delta(A) := \{a\in A, f^{-1}(a) \text{ is not smooth}\}$. For $\forall s\in A-\Delta(A)$, easy to compute  $\dim \pi_f^{-1}(s) = \dim X$. So we only need to prove for $\forall a \in \Delta(A)$, $\dim \pi_f^{-1}(a) = \dim X$. Fix $a_0 \in \Delta(A)$. Denote $\text{sing}(F_{a_0})$ the set of singular point of $F_{a_0}$. $F_{a_0}$ is integral, So $F_{a_0} - \text{sing}(F_{a_0})$ is  irreducible. Denote $\overline{T_f^* X|_{F_{a_0} - \text{sing}(F_{a_0})}}^{T^*X}$ the closure of  $T_f^* X|_{F_{a_0} - \text{sing}(F_{a_0})} $ in $T^*X$.
  %   We claim  $\pi_f^{-1}(a_0) = \overline{T_f^* X|_{F_{a_0} - \text{sing}(F_{a_0})}}^{T^*X}$. Granting this claim for the moment. 
   %  since $T_f^* X|_{F_{a_0} - \text{sing}(F_{a_0})} \subseteq \pi_f^{-1}(a_0)$, and $\pi_f^{-1}(a_0)$ is a closed subset in $T^*X$, we have  $\pi_f^{-1}(a_0) = \overline{T_f^* X|_{F_{a_0} - \text{sing}(F_{a_0})}}^{T^*X}  = \overline{T_f^* X|_{F_{a_0} - \text{sing}(F_{a_0})}}^{\pi_f^{-1}(a_0)}$. Note that $T_f^* X|_{F_{a_0} - \text{sing}(F_{a_0})}$ is irreducible since it's the total space of a vector bundle over an irreducible base space, then we have $\dim \pi_f^{-1}(a) = \dim  T_f^* X|_{F_{a_0} - \text{sing}(F_{a_0})} $. In fact $ \dim  T_f^* X|_{F_{a_0} - \text{sing}(F_{a_0})} = \dim X$, we will show it after we prove the claim. 

     
  %   As for the claim, fix $x_0\in F_{a_0}$, we only need to prove for any Caucahy sequence $\{(x_n,\eta_n)\ |\    (x_n,\eta_n)\in T_f^*X|_{X^\circ} \}$ with limit point $(x_0, \eta_0)\in T^*X$, there exists a Caucahy sequence $\{(y_n,\mu_n)\ |\ (y_n,\mu_n)\in T_f^*X|_{X^\circ},y_n\in F_{a_0} \}$ which has the same limit point $(x_0, \eta_0)\in T^*X$.

  %   Assume $\dim A = g$, take $e_1,...,e_g$ to be the basis of $T_{a_0}^* A$. Since A is abelian, we have a conanical isomorphism $H^0(A,\Omega_A^1)\simeq T_{a_0}^* A$. Denote by $\omega_1,...,\omega_g\in H^0(A,\Omega_A^1)$ the correspondence of $e_1,...,e_g$ respectively. For $\forall a\in A$, $\omega_1(a),...,\omega_g(a)$ forms a basis of $T^*_aA$. 

     
  %   For $\forall x\in X_0$, $d_{f,x}\colon T_xX\to T_{f(x)}A$ is surjective, then the dual map $d_{f,x}^*\colon T_{f(x)}^*A\to T_x^*X$ is injective. So $d_{f,x}^*T_{f(x)}^*A = \text{span}< f^*(\omega_1(f(x))),...,f^*(\omega_g(f(x)))> = \text{span}<f^*\omega_1(x),...,f^*\omega_g(x)>$.
     
  % For $\forall \eta\in d_{f,x}^*T_{f(x)}^*A$, we have $\eta(T_xF_{f(x)}) = 0$. It implies $d_{f,x}^*T_{f(x)}^*A \subseteq \phi_f^{-1}(x)$. Moreover, $\dim \phi_f^{-1}(x) = \dim X - \dim F_{f(x)}$. Since $f$ is flat, we know $\dim X - \dim F_{f(x)} = \dim A = g$. So in fact  $d_{f,x}^*T_{f(x)}^*A = \phi_f^{-1}(x)$.

    
    
  %   Given Cauchy sequence $\{(x_n,\eta_n)\ |\ (x_n,\eta_n)\in T_f^*X|_{X^\circ} \}$ with limit point $(x_0, \eta_0)\in T^*X$ . Since $x_n\in X_0$, we have $\eta_n \in \text{span}<f^*\omega_1(x_n),...,f^*\omega_g(x_n)>$. Denote $\eta_n = \sum_{i=1}^g \alpha_i^n f^*\omega_i(x_n)$. Define sequence $\{(y_n,\mu_n)\}$ as follows: Choose $y_n$ to be any sequence converge to $x_0$ such that $y_n\in F_{a_0}\cap X^\circ$. Define $\mu_n = \sum_{i=1}^g \alpha_i^n f^*\omega_i(y_n) $, clearly $\mu_n\in d_{f,y_n}^*T_{f(y_n)}^*A = \phi_f^{-1}(y_n)$, thus  $ (y_n,\mu_n)\in T_f^*X|_{X^\circ}\text{ and }y_n\in F_{a_0} $ for any $n$. By the continuity of $f^*\omega_i,i=1,...,g$, we conclude that the sequence $\{(y_n,\mu_n)\}$ is indeed a Cauchy sequence and has the same limit as $\{(x_n,\eta_n)\}$.

 %    At last, we prove $ \dim  T_f^* X|_{F_{a_0} - \text{sing}(F_{a_0})} = \dim X$. In fact, by the proof of the claim, the fiber of $T_f^* X|_{F_{a_0} - \text{sing}(F_{a_0})}\to F_{a_0} - \text{sing}(F_{a_0})$ has dimension $g$, so 
  %   $ \dim  T_f^* X|_{F_{a_0} - \text{sing}(F_{a_0})} = \dim (F_{a_0} - \text{sing}(F_{a_0}))+g = \dim F_{a_0} + \dim A$. Since $f$ is flat, we have $\dim F_{a_0} = \dim X- \dim A$.  So $ \dim  T_f^* X|_{F_{a_0} - \text{sing}(F_{a_0})} = \dim X$ as desired.
%\end{proof}

Next we have the following two lemmas, which will be used in the proof of Theorem \ref{thm:simple-smooth}.

\begin{lemma}[{\cite[Proposition 3.2]{DHL24}}]\label{prop:DHL}
   Let $A$ be an abelian variety and $Z$ be a proper irreducible subvariety of $A$. Then the following are equivalent
   
a) $Z$ is not fibred by tori and $\dim Z > 0$;

b) General holomorphic 1-form $\omega\in H^0(A,\Omega^1_A)$ restricted to $Z^{\text{reg}}$, i.e., $\omega|_{Z^{\text{reg}}}$ admits isolated
zeros on the smooth locus $Z^{\text{reg}}$. 

 In particular, let $B \subseteq A$ be the largest (in dimensional sense) Abelian subvariety such that $Z$ is fibred by $B$, we have $p_2(T_Z^*A) = H^0(A/B,\Omega^1_{A/B})$, where $T_Z^*A$ is the conormal variety\footnote{For a subvariety $Z\subset A$, the conormal variety $T^*_{Z}A$ is defined to be the closure of the conormal bundle of the smooth locus of $Z$ in $T^*A$.} of $Z$ in $A$, and $p_2\colon T^*A\to H^0(A, \Omega_A^1)$ is the natural projection.
\end{lemma}

\begin{lemma}\label{lem:smooth-all-forms}
   Let $f\colon X\to A$ be a morphism from a smooth projective variety $X$ to an abelian variety $A$. Then $f$ is smooth, if and only if for any nonzero holomorphic 1-form $\omega\in H^0(A, \Omega_A^1)$, $f^*\omega$ has no zero.
\end{lemma}

\begin{proof}
$f$ is smooth if and only if $f$ is a submersion, i.e., the tangent map of $f$ is surjective at any closed point of $X$. This is equivalent to that $f^*\omega$ has no zero for any nonzero holomorphic 1-form $\omega\in H^0(A, \Omega_A^1)$.
\end{proof}

Now we use Theorem \ref{thm:homological-trivial-deformation-no-blow-up-codim0}, Lemma \ref{prop:DHL} and Lemma \ref{lem:smooth-all-forms} to show the following

\begin{theorem}\label{thm:simple-nowhere vanishing 1-form}
    Let $f\colon X\to A$ be a morphism from a smooth projective variety $X$ to a simple abelian variety $A$. The following are equivalent:

    1) $f$ is smooth;
    
    2) there exists a holomorphic 1-form $\omega\in H^0(A, \Omega^1_A)$ such that $f^*\omega$ has no zero.
\end{theorem}



\begin{proof}
1) implies 2) is trivial. We show 2) implies 1). For the projective morphism $f\colon X\to A$, we may choose an embedding $i\colon X\to \mathbb{P}^m\times A$ for some integer $m$ such that $f$ and $i$ fit into the following commutative diagram.
\[\xymatrix{
X\ar[rd]_-f\ar@{^{(}->}[r]^-i&\mathbb{P}^m\times A  \ar[d]^-q \\
&A,
}\]
where $q$ is the natural projection. Then we have the commutative diagram 
\[\xymatrix{
T^*_X(\mathbb{P}^m\times A) \ar@{^{(}->}[r]&T^*(\mathbb{P}^m\times A)\\
f^*T^*A  \ar@{^{(}->}[r]^-{i\times \id}\ar[d]&q^*T^*A\ar[d] \ar@{^{(}->}[u]\ar[r]^-{q\times \id}&T^*A\ar[d]\ar[r]^-{p_2}& H^0(A, \Omega_A^1)  \\
X\ar[r]^-{i}&\mathbb{P}^m\times A\ar[r]^-q &A,
}\]
We claim that $$p_2\circ(q\times \id)(q^*T^*A\cap T_X^*(\mathbb{P}^m\times A))=\{\omega\in H^0(A, \Omega_A^1)\ |\ Z(f^*\omega)\neq \emptyset\},$$
where $Z(f^*\omega)$ is the zero loci of holomorphic 1-form $f^*\omega$. In fact, it is clear that \begin{align*}
   (q\times \id)(q^*T^*A\cap T_X^*(\mathbb{P}^m\times A))=&(f\times \id)(f^*T^*A\cap T_X^*(\mathbb{P}^m\times A)) \\
   =&(f\times \id)((X\times H^0(A, \Omega_A^1))\cap T_X^*(\mathbb{P}^m\times A)).
\end{align*} 
Then by the exact conormal sequence of vector bundles $$0\to T^*_X(\mathbb{P}^m\times A)\to T^*(\mathbb{P}^m\times A)|_X\to T^*X\to 0,$$ we infer that the above claim holds. 

For the rest, we show that for  any nonzero holomorphic 1-form $\omega$ on $A$, $f^*\omega$ has no zero. By assumption 2), one has that the underlying differentiable structure of $X$ is a differentiable fibre bundle over the circle $S^1$ by \cite[Proposition 7]{Kot22}. Then by \cite[Corollary 3.4]{SY25} or \cite[Corollary 1.5]{DHL24}, we have $R^if_*\mathbb{Z}$ are all $\mathbb{Z}$-local systems for all integers $i$. Now we apply Theorem \ref{thm:homological-trivial-deformation-no-blow-up-codim0} and get that $f$ is without blow-up in codimension 0. Therefore, by \cite[Proposition 3.3, 2)]{Sab85}, $(q\times \id)(q^*T^*A\cap T_X^*(\mathbb{P}^m\times A))$ is a Lagrangian subvariety of $T^*A$, i.e. it is a finite union of conormal varieties $T^*_{Z_i}A$ along various subvarieties $Z_i$ of A. Since $A$ is simple, any positive dimensional proper subvariety of $A$ is of general type. Then by Lemma \ref{prop:DHL}, we have $$p_2\circ(q\times \id)(q^*T^*A\cap T_X^*(\mathbb{P}^m\times A))=\{0\}\ \text{or}\ H^0(A,\Omega_A^1),$$ so is $\{\omega\in H^0(A, \Omega_A^1)\ |\ Z(f^*\omega)\neq \emptyset\}$. However, we have a nowhere vanishing holomorphic 1-form $f^*\omega$. Hence $$\{\omega\in H^0(A, \Omega_A^1)\ |\ Z(f^*\omega)\neq \emptyset\}=\{0\},$$ which means for any nonzero holomorphic 1-form $\omega$ on $A$, $f^*\omega$ has no zero. Thus $f$ is smooth by Lemma \ref{lem:smooth-all-forms}.
 \end{proof}




\begin{remark}
    In Theorem \ref{thm:simple-smooth}, one cannot expect that $f$ is an analytic fibre bundle or isotrivial. For example, let $Y$ be a smooth projective variety of general type, which contains an elliptic curve $E$. Denote by $\sigma$ the closed embedding $E\to Y$, and consider the trivial family $p_1\colon E\times Y\to E$. Let $\Gamma_{\sigma}$ be the graph of $\sigma$ in $E\times Y$. Denote the blowup of $E\times Y$ along $\Gamma_{\sigma}$ by $X$. Then the morphism $X\to E$ is smooth, but not isotrivial.  
\end{remark}


\section{Linearity of the space of holomorphic 1-forms with zeros}\label{sect:deg-1-form-linear}
%\section{Criterion for detecting Hodge one-forms}
Let $X$ be a smooth projective variety. For any holomorphic 1-form $\omega\in H^0(X, \Omega_X^1)$, $Z(\omega)$ is denoted to be the set-theoretic zero loci of $\omega$ in $X$. We recall the following algebraic invariant introduced by Budur-Wang-Yoon \cite{BWY16}. 

\begin{definition}
The \textit{$i$-th lower level set} of $X$ is defined to be 
\[
W^i(X):=\{\omega\in H^0(X, \Omega_X^1)\ |\ \codim Z(\omega)\leq i \}.
\] 
\end{definition}

For a smooth projective variety $X$ of dimension $n$, $W^n(X)$ is the whole space of holomorphic 1-forms with zeros on $X$. Inside of $W^i(X)$, one has subsets arising from the tangent cone of the cohomology jump loci, which is very well studied,
\[
V^i(X):=\{\omega\in H^0(X, \Omega_X^1)\ |\ (H^{\bullet}(X,\CC),\wedge\omega)\text{ is not exact at } H^j(X,\CC) \text{ for } j\leq i\}.
\]
Consider the character variety $\Hom(\pi_1(X),\CC^{\ast})$ of $X$ which is a direct product of $(\CC^{\ast})^{b_1(X)}$ with a finite group. Let $\cV^i(X)$ be the subset of $\Hom(\pi_1(X),\CC^{\ast})$ consisting of rank one local systems $\cL$ on $X$ such that $H^{j}(X,\cL)\neq 0$ for $j\leq i$. By the tangent cone theorem \cite[Theorem A]{DPS09}, $V^i(X)$ is the intersection of $H^0(X, \Omega_X^1)$ with the tangent cone of $\cV^i(X)$ at the trivial local system. Therefore, $V^i(X)$ is \textit{linear}, i.e., a union of finitely many vector subspaces in $H^0(X, \Omega_X^1)$ (See \cite{Bea88} and \cite{Sim93}).


One can show that $V^i(X)\subseteq W^i(X)$. In fact, suppose by contradiction that there is a holomorphic 1-form $\omega\in V^i(X)$ and $\omega\notin W^i(X)$. Then $\codim Z(\omega)>i$. By \cite[Proposition~3.4]{GL87}, $(H^{\bullet}(X,\CC),\wedge\omega)$ is exact at any degree $\leq i$, which is a contradiction.  
%Observe that for any codimension $k$ degenerate 1-form $\omega$ on $X$, the sequence 
%\[
 %\rightarrow H^{i-1}(X,\CC) \xrightarrow{\wedge \omega} H^{i}(X,\CC) \xrightarrow{\wedge \omega} H^{i+1}(X,\CC)\xrightarrow{\wedge \omega} \cdots 
 %\]
 %is exact at $H^i(X,\CC)$ for $i<k$, in particular, $\omega\notin V^{k-1}$. If $\codim Z(\omega)=i$ and $\omega\notin V^k$, then above sequence is exact for all $i\leq k$.


     
Motivated by the linearity of $V^n(X)$, it is natural to ask whether $W^n(X)$ is linear or not \cite[Question 1.10]{DHL24}. There are some known cases. For example,  if $X$ is of general type or has nonzero topological/holomorphic Euler characteristic, then $W^n(X)$ is linear (see \cite{PS14} and \cite{GL87}). %For smooth projective varieties of $n\leq 3$, $W^n(X)$ is linear by \cite{HS21}.
If the Albanese map of $X$ is finite, then $W^n(X)$ is linear by \cite{DHL24}. Moreover, if $X$ is a subvariety of an abelian variety, then $W^n(X)$ is linear by \cite{Ueno75}. Based on works of \cite{Chu24} and \cite{HWZ26}, we obtain the linearity property for any smooth good minimal models.

\begin{theorem}\label{thm:linearity-zeros-gmm}
    If $X$ is a smooth projective good minimal model (i.e., $K_{X}$ is semiample) of dimension $n$, then $$W^{n}(X)=\{\omega\in H^{0}(X,\Omega_{X}^{1})~|~Z(\omega)\neq \emptyset\}$$ is a  vector subspace of $H^{0}(X,\Omega_{X}^{1})$. 
\end{theorem}
\begin{proof}
    If every holomorphic 1-form on $X$ has zeros, then $W^{n}(X)=H^{0}(X,\Omega_{X}^{1})$. Now we may assume that $X$ admits a nowhere vanishing holomorphic 1-form. Then by \cite{Chu24} or \cite{HWZ26}, there exists an isotrivial smooth fibration $f\colon X\to A$ where $A$ is a positive dimensional abelian variety. Moreover, there exists a finite 
    \'etale cover $A'\to A$ such that the base change of $X$ along $A'\to A$ splits, that is, $X':=X\times_{A}A'\cong Y\times A'$, where $Y$ is a smooth projective variety. Since $K_{X'}$ is semiample, it follows that $K_{Y}$ is semiample. 
    Iterating the above process, we can assume that there exists a finite \'etale cover $\pi:\tilde{X}\to X$ such that $\tilde{X}\cong \tilde{Y}\times B$, where $\tilde{Y}$ is a smooth projective variety without nowhere vanishing holomorphic 1-forms, and $B$ is a positive dimensional abelian variety. Denote by $p_{1}:\tilde{X}\to \tilde{Y}$ the projection to the first factor. Since any nonzero holomorphic 1-form on $B$ has no zeros, we have $W^{n}(\tilde{X})=p_{1}^{*}H^{0}(\tilde{Y},\Omega_{\tilde{Y}}^{1})$, which is a linear subspace of $H^{0}(\tilde{X},\Omega_{\tilde{X}}^{1})$. To show that $W^{n}(X)\subset H^{0}(X,\Omega_{X}^{1})$ is a linear subspace, it suffices to verify that $W^{n}(X)$ is closed under addition. Since $\pi$ is \'etale, for any two holomorphic 1-forms $\omega_{1},\omega_{2}\in W^{n}(X)$, $\pi^{*}(\omega_{1}),\pi^{*}(\omega_{2})\in W^{n}(\tilde{X})$. As $W^{n}(\tilde{X})$ is a linear subspace, we have $\pi^{*}(\omega_{1})+\pi^{*}(\omega_{2})=\pi^{*}(\omega_{1}+\omega_{2})\in W^{n}(\tilde{X})$. Using again that $\pi$ is \'etale we conclude $\omega_{1}+\omega_{2}\in W^{n}(X)$.   
\end{proof}

For varieties $X$ with simple Albanese varieties, we obtain the linearity property of $X$ by Theorem \ref{thm:simple-smooth}.

\begin{theorem}\label{prop:simple-linear}
  Let $f\colon X\to A$ be a morphism from an $n$-dimensional smooth projective variety $X$ to a simple abelian variety $A$,  then $W^n(X)\cap f^*H^0(A,\Omega_A^1)=f^*H^0(A,\Omega_A^1)$ or $\{0\}$. In particular, if the Albanese variety of $X$ is simple, then $W^n(X)=H^0(X, \Omega_X^1)$ or $\{0\}$. 
\end{theorem}

\begin{proof}
This trivially follows from Theorem \ref{thm:simple-smooth}.
\end{proof}

In general, for an $n$-dimensional smooth projective variety $X$, $W^n(X)$ is not always a vector subspace.
\begin{example}
    Let $Y=E_{1}\times E_{2}\times Z$, where $E_{1}$, $E_{2}$ are two non-isogenous elliptic curves, $Z$ is a simply connected smooth projective variety of dimension $n-2$. Let $p,q\in Z$ be two distinct closed points, and let $C_{1}=E_{1}\times\{e_{2}\}\times\{p\}$, $C_{2}=\{e_{1}\}\times E_{2}\times\{q\}$. Let $\rho:X\to Y$ be the blow-up of $Y$ along $C_{1}\cup C_{2}$. Then we have
    $$W^{n}(X)=p_{1}^{*}H^{0}(X,\Omega_{E_{1}}^{1})\cup p_{2}^{*}H^{0}(X,\Omega_{E_{2}}^{1}),$$
    where $p_{i}\colon X\to E_{i}$ $(i=1,2)$ are the natural projections. In particular, $W^{n}(X)$ is a linear subset of $H^{0}(X,\Omega_{X}^{1})$. 
\end{example}

 

\begin{remark}
In general, the structure of $W^n(X)$ is very sensitive to the birational modifications involved in the minimal model program from a smooth projective variety to its minimal models. 
\end{remark}

For the remaining part of this section, we discuss the linearity of $W^n(X)$ for $n=\dim X\leq 3$.

\begin{lemma}\label{lem:linearity-blowup}
    Let $X$ be an n-dimensional smooth projective variety and $C\subset X$ be a smooth projective subvariety of dimension $d$. Let $Y=\Bl_{C}X$ be the blow-up of $X$ along $C$. If $W^n(X)$ and $W^d(C)$ are both linear, then so is $W^n(Y)$.
\end{lemma}

\begin{proof}
  For the blow-up morphism $\pi\colon Y\to X$, we have the following two simple observations. We leave the proof to the reader. 
  \begin{itemize}
        \item[(a)] the set $\{\eta\in H^0(X, \Omega_X^1)|Z(\eta|_{C})\neq \emptyset\}$ is a linear subspace of $H^0(X, \Omega_X^1)$;
        \item[(b)] for $\pi^{\ast}\omega\in H^0(Y, \Omega_Y^1)$, $\pi^{\ast}\omega$ has zeros along $\pi^{-1}(C)$ if and only if  $Z(\omega|_C)\neq\emptyset$.
    \end{itemize} 
Then by these observations, we have  
\begin{align*}
W^n(Y) =&\{\pi^{\ast}\omega \in \pi^{\ast} H^0(X, \Omega_X^1)\ |\ Z(\pi^{\ast}\omega)\cap (Y-\pi^{-1}(C))\neq\emptyset\}\\
     &\cup \{\pi^{\ast}\omega \in \pi^{\ast} H^0(X, \Omega_X^1)\ |\ Z(\pi^{\ast}\omega)\cap \pi^{-1}(C)\neq\emptyset\} \\
     =& \pi^{\ast} \{\omega\ |\ Z(\omega)\cap (X-C)\neq\emptyset\}\cup  \pi^{\ast} \{\omega\ |\ Z(\omega|_C)\neq\emptyset\}  \\
     =& \pi^{\ast} W^n(X)\cup \pi^{\ast}\{\omega\ |\ Z(\omega|_C)\neq\emptyset\}.
\end{align*}
Hence $W^n(Y)$ is linear. 
\end{proof}




\begin{theorem}\label{prop:23linearity}
 Let $X$ be a smooth projective variety with $\dim X\leq 3$, then $W^{\dim X}(X)$ is linear.   
\end{theorem}

\begin{proof}
    This follows directly from Theorem \ref{thm:linearity-zeros-gmm}, Lemma \ref{lem:linearity-blowup}, \cite[Corollary 3.1]{Sch21}, and \cite[Theorem 1.3]{HS21}.
\end{proof}










\section{Non-linearity of lower level sets and Non-formal 1-forms}
In this section, we give an example for a smooth projective subvariety of an abelian variety with a nonlinear lower level set. Moreover, we study the related holomorphic 1-forms that have zeros and do not come from cohomology jump loci, and smooth projective surfaces with this kind of holomorphic 1-forms.

\subsection{Nonlinear Examples}


\begin{theorem}\label{thm:non-linear example}
There exists a $4$-dimensional smooth projective subvariety $X$ of an abelian 6-fold, such that the third lower level set $W^3(X)$ is not linear.
\end{theorem}
 
To deal with the above theorem, we first introduce the following lemma, which is known to experts. The proof presented here is due to Christian Schnell.

\begin{lemma}\label{lem:conormalspace-nonlinear}
  Let $X$ be a smooth projective variety, and let $f\colon X\to A$ be a morphism to an abelian variety $A$. Denote $$Z_{f}=\{(x,\omega)\in X\times H^{0}(A,\Omega^{1}_{A}) ~|~~ f^{*}\omega(x)=0 \}$$ and let $$S_{f}\coloneqq(f\times \mathrm{id})(Z_{f})=\{(a,\omega)\in A\times H^{0}(A,\Omega^{1}_{A}) ~|~~  f^{*}\omega(x)=0,\ \text{for some }x\in f^{-1}(a) \},$$ then we have 
$\mathrm{dim} S_{f}\leq\mathrm{dim} A$, and every irreducible component of $S_{f}$ of dimension equal to $\mathrm{dim}A$ is a conormal variety.   
\end{lemma} 

\begin{proof}
Let $W\subseteq S_{f}$ be an irreducible component. Denote by $p_{1}\colon S_{f}\to A$, $p_{2}\colon S_{f}\to H^{0}(A,\Omega^{1}_{A})$ the two natural projections. Let $Z=p_{1}(W)\subseteq A$. Note that $S_{f}$ is conical and $W$ is closed in $ A\times H^{0}(A,\Omega^{1}_{A})$, it follows that $Z$ is closed in $A$. It suffices to show that $W$ is contained in the conormal variety $ T^{*}_{Z}A$. By definition, for every pair $(a,\omega)\in W$, there is a point $x\in f^{-1}(a)$ such that $\omega$ vanishes on the image of the tangent map $f_{*,x}\colon T_{x}{X}\to T_{a}A$. Claim that for a general closed point $a \in Z^{reg}$, we have $$f_{*,x}(T_{x}{X})\supseteq T_{a}Z.$$
In fact, we note that $Z\subseteq f(X)$. Take a resolution of singularities $\eta\colon \widetilde{f^{-1}(Z)}\to f^{-1}(Z)$ of the algebraic set $f^{-1}(Z)$. By generic smoothness, one can choose a nonempty open subset $U$ in $Z$ such that $f\circ \eta\colon (f\circ \eta)^{-1}(U)\to U$ is smooth. Denote $V\coloneqq (f\circ \eta)^{-1}(U)$. For any $a\in U$, and any $x\in f^{-1}(a)$, we take $\tilde{x}\in V$ such that $\eta(\tilde{x})=x$. Denote the natural map $V\to X$ by $\varphi$. Then we have $$T_aZ=(f\circ\eta)_{*,\tilde{x}}(T_{\tilde{x}}V)=f_{*,x}(\varphi_{*,\tilde{x}}(T_{\tilde{x}}V))\subseteq f_{*,x}(T_{x}{X}).$$
It follows that $(a,\omega)\in T^{*}_{Z^{reg}}A$, hence $W \subseteq T^{*}_{Z}A$.
\end{proof}


\begin{proof}[Proof of Theorem \ref{thm:non-linear example}]
Take any simple abelian $5$-fold $A$ and $Y\subset A$ a general hyperplane section, in particular, $Y$ is a smooth projective $4$-fold of general type, and the normal bundle $\cN_{Y/A}$ of $Y$ in $A$ is ample, since $\cN_{Y/A}=\cO_{A}(Y)|_{Y}$. Consider the conormal sequence \[0\to \cN^*_{Y/A}\to Y\times H^0(A, \Omega_A^1)\to \Omega_Y^1\to 0\] and the short exact sequence of fibres over each closed point $y\in Y$ 
\begin{equation}\label{eq:conormal-seq}
  0\to \cN^*_{Y/A, y}\to H^0(A, \Omega_A^1)\to \Omega^1_{Y,y}\to 0  
\end{equation}
 We then have the induced Gauss map  from $Y$ to the Grassmannian \[
 \cG\colon Y\to \textnormal{Gr}(1,H^0(A,\Omega_A^1)); y\mapsto [\cN^*_{Y/A, y}],
 \]
By \cite[Corollary~2]{Ran84} that $\cG$ is a finite morphism onto its image. Moreover, by \cite[Theorem]{WZ25}, one can pick a general irreducible smooth projective  curve $C\subset Y$ with the property that $\cG(C)\subset \textnormal{Gr}(1,H^0(A,\Omega_A^1))\cong \PP^4$ is not contained in any hyperplane of $\textnormal{Gr}(1,H^0(A,\Omega_A^1))\cong \PP^4$ (i.e., nondegenerate), and the set \[\cS_1:=\{\omega\in H^0(A,\Omega_A^1)\ |\ \omega|_Y \text{ has isolated zeros on }C\}\cup\{0\}\] is non-linear. For reader's convenience, we explain the non-linearity of $\cS_1$ here. In fact, by the conormal sequence (\ref{eq:conormal-seq}), $$\cS_1=\bigcup_{y\in C}\cN^*_{Y/A, y}=\bigcup_{[\cN^*_{Y/A, y}]\in \mathcal{G}(C)}\cN^*_{Y/A, y},$$ which is not linear, since $\mathcal{G}(C)$ is not contained in any hyperplane. Moreover, we observe that $\cS_1$ is irreducible and not contained in any proper vector subspace of $H^0(A,\Omega_A^1)$.


Now take any elliptic curve $E$ and consider the abelian 6-fold $E\times A$. Inside $E\times A$, we have smooth subvarieties $E\times C\subset E\times Y\subset E\times A.$ We will show that one can choose a smooth hyperplane section $X$ of $E\times Y$ that passes through $E\times C$, and the third lower level set $W^3(X)$ is not linear. We will proceed with several steps.

\textbf{Step 1. We show that there exists a smooth hyperplane section $X$ of $E\times Y$ such that $E\times C\subset X$.}

In fact, fix $L$ a very ample divisor on $E\times Y$, and denote by $\mathcal{I}_{E\times C}$ the ideal sheaf of $E\times C$. We  consider the short exact sequence %(after twisted by $L^{\otimes m}$ for any integer $m>0$)
\[
0\to \cI_{E\times C}\otimes L^{\otimes m}\to\cO_{E\times Y}\otimes L^{\otimes m}\to\cO_{E\times C}\otimes L^{\otimes m}\to 0
\]
and the induced long exact sequence
\[
0\to H^0(\cI_{E\times C}\otimes L^{\otimes m})\to H^0(\cO_{E\times Y}\otimes L^{\otimes m})\to H^0(\cO_{E\times C}\otimes L^{\otimes m})\to H^1(\cI_{E\times C}\otimes L^{\otimes m})
\]
By asymptotic Riemann-Roch and Serre vanishing, the leading term of $h^0(\cI_{E\times C}\otimes L^{\otimes m})$ is $am^5$ for some $a>0$ and $H^1(\cI_{E\times C}\otimes L^{\otimes m})=0$  for sufficiently large $m$. Then for $m\gg 0$, inside of the complete linear system $|mL|$, we have the linear system $\mathfrak{d}$ corresponding to $H^0(\cI_{E\times C}\otimes L^{\otimes m})$ whose scheme-theoretic base locus is $E\times C$. Then by the strong Bertini theorem \cite[Theorem 3.1]{DH91}, since $2\dim (E\times C)\leq \dim (E\times Y)$, there exists a hyperplane section $X$ in $\mathfrak{d}$ such that $X$ is smooth and contains $E\times C$.  


\textbf{Step 2. We may assume that for the chosen hyperplane section $X$ in Step 1, the natural map $\pi\colon X\hookrightarrow E\times Y\to Y$ is finite away from $E\times C$.}

We take the embedding $E\times Y\subset \mathbb{P}^n$ given by $|mL|$. For a general hyperplane section $H$ of $E\times Y$, $H$ intersects curve $E_y\coloneqq E\times \{y\}$ for any closed point $y\in Y$. Hence the natural projection $H\hookrightarrow E\times Y\to Y$ is surjective and generically finite. Denote the inclusion $E\times Y\subset \PP^n$ by $i$. Then the following morphism $$\varphi\colon E\times Y\to \PP^n\times Y; (e,y)\mapsto (i(e,y), y)$$ gives a bounded smooth family of elliptic curves in $\PP^n$ over $Y$. Then we may assume that for the chosen $m$ in Step 1, and each closed point $y\in Y$,
\begin{equation}\label{eq:vanishingH1}
 H^1(E\times Y,  \cI_{E_y}\otimes L^{\otimes m})=H^1(E\times Y,  \cI_{(E\times C)\cup E_y}\otimes L^{\otimes m})=0,   
\end{equation}
 and 
 \begin{equation}\label{eq:H0>4}
  \dim H^0(E_y, L^{\otimes m}|_{E_y})\gg \dim Y=4.   
 \end{equation}
For an arbitrary closed point $y\in Y-C$, consider the inclusion $$E\times C\hookrightarrow (E\times C)\cup E_y,$$ where $(E\times C)\cup E_y$ is a disjoint union. Then we have a commutative diagram
$$\xymatrix{
	0\ar[r] & \mathcal{I}_{(E\times C)\cup E_y}\otimes L^{\otimes m} \ar@{^{(}->}[d] \ar[r] &L^{\otimes m} \ar@{=}[d]\ar[r]&L^{\otimes m}|_{(E\times C)\cup E_y} \ar@{->>}[d] \ar[r]&0 \\
0\ar[r] & \mathcal{I}_{E\times C}\otimes L^{\otimes m}   \ar[r] &L^{\otimes m}  \ar[r]&L^{\otimes m}|_{E\times C} \ar[r]&0. 
}$$

Hence, by Equation (\ref{eq:vanishingH1}), we get the following commutative diagram 
$$\xymatrix{
	0\ar[r] & H^0(\mathcal{I}_{(E\times C)\cup E_y}\otimes L^{\otimes m}) \ar@{^{(}->}[d] \ar[r] &H^0(L^{\otimes m}) \ar@{=}[d]\ar[r]& H^0(L^{\otimes m}|_{(E\times C)\cup E_y}) \ar@{->>}[d] \ar[r]&0 \\
0\ar[r] & H^0(\mathcal{I}_{E\times C}\otimes L^{\otimes m})   \ar[r] & H^0(L^{\otimes m})  \ar[r]& H^0(L^{\otimes m}|_{E\times C}) \ar[r]&0. 
}$$

Hence \begin{align*}
	\dim H^0(\mathcal{I}_{E\times C}\otimes L^{\otimes m})-\dim H^0(\mathcal{I}_{(E\times C)\cup E_y}\otimes L^{\otimes m}) \\     
	  = \dim H^0(L^{\otimes m}|_{(E\times C)\cup E_y})-\dim H^0(L^{\otimes m}|_{E\times C}) \\
	  =\dim H^0(L^{\otimes m}|_{E_y})\\
	  \gg 4,
\end{align*}
by Equation (\ref{eq:H0>4}).

Consider the dual space $\PP^{n*}$ of hyperplanes in $\PP^n$, and the following incidence set $$A\coloneqq \{(H, y)\in \PP^{n*}\times (Y-C)\ |\ E_y\subset H\},$$ together with two natural projections $p_1\colon A\to \PP^{n*}$ and $p_2\colon A\to Y-C$. We need to show that the intersection $p_1(A) \cap \mathbb{P}H^0(\mathcal{I}_{E\times C}\otimes L^{\otimes m})$ is a countable union of proper subschemes of $\mathbb{P}H^0(\mathcal{I}_{E\times C}\otimes L^{\otimes m})$. It suffices to show that
$$\dim p_1(A)\cap \PP H^0(\mathcal{I}_{E\times C}\otimes L^{\otimes m})< \dim \PP H^0(\mathcal{I}_{E\times C}\otimes L^{\otimes m})$$ since $A$ has a scheme structure. Consider the following diagram
\begin{equation}\label{dia:full-construction}
\xymatrix{
	&&A \ar[ld]_-{p_1}\ar[rd]^-{p_2}   \\
\PP H^0(\mathcal{I}_{E\times C}\otimes L^{\otimes m})\ar@{^{(}->}[r]	&\PP^{n*} && Y-C.
}
\end{equation}
For each closed point $y\in Y-C$, the fibre $p_2^{-1}(y)$ is a linear subspace of $\PP^{n*}\times \{y\}$.  Also, \begin{align*}
\dim p_1p_2^{-1}(y)\cap \PP H^0(\mathcal{I}_{E\times C}\otimes L^{\otimes m})\\
=\dim \PP H^0(\mathcal{I}_{(E\times C)\cup E_y}\otimes L^{\otimes m})\\
\ll \dim \PP H^0(\mathcal{I}_{E\times C}\otimes L^{\otimes m})-4
\end{align*} Hence $$\dim p_1(A)\cap \PP H^0(\mathcal{I}_{E\times C}\otimes L^{\otimes m})\ll \dim \PP H^0(\mathcal{I}_{E\times C}\otimes L^{\otimes m}).$$ 
\vspace{0.03em}


%Now we show that $A$ is a  subscheme of $\PP^{n*}\times (Y-C)$. Since $E\times (Y-C)\to Y-C$ is a flat family of curves in $\PP^n$, $E_y$ have the same Hilbert polynomial, denoted as $\Phi$, for all $y\in Y$. Now consider the Hilbert scheme $\text{Hilb}_{\PP^n}^{\Phi}$ of curves in $\PP^n$ with Hilbert polynomial $\Phi$, together with the universal family $\cE\to \text{Hilb}_{\PP^n}^{\Phi}$. By the universal property, the flat family $E\times (Y-C)\to Y-C$ is induced from a morphism $f\colon Y-C \to \text{Hilb}_{\PP^n}^{\Phi}$. Note that the scheme-theoretic intersection $H\cap E=E$ gives rise to a subscheme structure $$\cA\coloneqq \{(H, [E])\in \PP^{n*}\times\text{Hilb}_{\PP^n}^{\Phi}  | E\subset H\}.$$ The pullback of $\cA$ to $\PP^{n*}\times Y-C$ via $f\colon Y-C\to \text{Hilb}_{\PP^n}^{\Phi}$ is $A$.

To formulate Step 3, we fix the following notations. By the Lefschetz hyperplane theorem, the inclusion $Y\subset A$ is the Albanese map of $Y$, and $X\subset E\times Y$ have the same Albanese. Hence $$H^0(X, \Omega_X^1)\cong H^0(E\times Y, \Omega_{E\times Y}^1)= q_1^*H^0(E, \Omega_E^1)\oplus q_2^*H^0(A, \Omega_A^1),$$ where $q_1\colon E\times Y\to E$ and $q_2\colon E\times Y\to Y\hookrightarrow A$ are natural maps.  


\textbf{Step 3. For any holomorphic 1-form $\omega=\omega_1+\omega_2$ on $E\times Y$  where $0\neq\omega_1\in q_1^*H^0(E, \Omega_E^1)$ and $\omega_2\in q_2^*H^0(A, \Omega_A^1)$, we show that $\omega|_X$ has only isolated zeros.}

Note that the normal bundle $$\cN_{X/E\times Y}\cong \cO_{E\times Y}(X)|_X=\cO(mL)|_X,$$ which is an ample line bundle on $X$. Also, $\omega=\omega_1+\omega_2$ is a holomorphic 1-form without zero on $E\times Y$ since $\omega_1\neq 0$.

Now suppose by contradiction that $\omega|_X$ admits a positive-dimensional zero loci $Z(\omega|_X)$, then there exists a projective integral curve $B$ contained in $Z(\omega|_X)$ set-theoretically. Consider the conormal sequence of $X$ inside $E\times Y$ 
	$$
\xymatrix{
	0 \ar[r] & \cN_{X/E\times Y}^* \ar[r] & \Omega^1_{E\times Y}|_{X} \ar[r]^-{\delta}&\Omega^1_{X} \ar[r] &0
}
$$ \[\delta\colon \omega\mapsto \omega|_X\] For any closed point $x\in B$, $\omega|_X(x)=0$ in the fibre $\Omega_{X,x}^1$, which means $\omega(x)$ is in the fibre $\cN_{X/E\times Y,x}^*$. Since $\omega$ is a holomorphic 1-form without zero on $E\times Y$, we have that $\omega$ gives rise to a nowhere vanishing section in $H^0(B, \cN_{X/E\times Y}^*|_{B})$. Taking the normalization $\xi\colon \widetilde{B}\to B$, we have $$\dim H^0(\widetilde{B}, \xi^*\cN_{X/E\times Y}^*|_{B})\geq \dim H^0(B, \cN_{X/E\times Y}^*|_{B})>0,$$
in particular, $\deg \xi^*\cN_{X/E\times Y}^*|_{B}>0$. Moreover, since $\cN_{X/E\times Y}$ is  an ample line bundle on $X$, $\xi^*\cN_{X/E\times Y}|_{B}$ is ample on $\widetilde{B}$. Then $\deg \xi^*\cN_{X/E\times Y}^*|_{B}<0$, which is a contradiction. Thus we are done for Step 3.
\vspace{0.5cm}

By Step 3, all holomorphic 1-forms $\eta\in W^3(X)$ are of the type $\eta=\omega|_X$, where  $\omega\in q_2^*H^0(A, \Omega_A^1)$, i.e., we have \begin{equation}\label{eq:w1-from-Y}
    W^3(X)\subseteq \pi^*H^0(Y, \Omega_Y^1)=\pi^*k^*H^0(A,\Omega_A^1),
\end{equation} where $k$ denotes the inclusion map $Y\to A$. Now we consider the following commutative diagram (\ref{dia:full-construction}) and fix several notations.
\begin{equation}\label{dia:full-construction}
\xymatrix{
	 E\times C \ar[r]^-{p'}\ar@{^{(}->}[d]_{\beta}  \ar@/_4 pc/[dd]_{\alpha} &C\ar@{^{(}->}[d]^{\gamma}\\
   E\times Y \ar[r]^-p & Y \ar@{_{(}->}[d]^k\\  
	X\ar@{->>}[ru]^{\pi}  \ar[r]_-f \ar@{_{(}->}[u]^j & A 
}
\end{equation}

By the purity theorem of  branched loci (\cite[Proposition~2]{Zar58}), the branched loci of the morphism $X\setminus(E\times C)\to Y\setminus C$ is a divisor $D$. Denote the image of $D$ in $Y\setminus C$ by $\Delta$. Then for any closed point $x\in D$, the tangent map $\pi_{\ast,x}$ at $x$ is not surjective, and  $\pi$ is \'etale away from $Y-(C \cup \Delta)$. By \cite[Proposition~6.3.10]{PAG2}, for any holomorphic 1-form $\mu$ on $Y$, $\mu$ has isolated zeros on $Y$. It follows that for any 
holomorphic 1-form $\eta\in W^3(X)$, $\eta|_{X-((E\times C)\cup D)}$ has at most isolated zeros.




\textbf{Step 4. For a general holomorphic 1-form $\xi$ on $A$, we show that $\pi^{\ast}(\xi|_Y)|_{X-(E\times C)}$ has only isolated zeros.}

Recall that the set $$\mathcal{S}_1:=\{\omega\in H^0(A,\Omega_A^1)\ |\ \omega|_Y \text{ has  isolated zeros on }C\}\cup\{0\}$$ is not linear, in particular, $\mathcal{S}_1\neq H^0(A,\Omega_A^1)$. Thus for a general holomorphic 1-form $\xi$ on $A$, $\xi|_Y$ has isolated zeros on $Y-C$ by \cite[Proposition~6.3.10]{PAG2}. Now suppose by contradiction that for a general holomorphic 1-form $\xi\in H^0(A, \Omega_A^1)$, $$\dim\pi^{\ast}(\xi|_Y)|_{X-(E\times C)}\geq 1.$$
 
 
 
 
 Consider the morphism $f=k\circ\pi\colon X \to A$, which is generically finite with exceptional set $E\times C$, and
 $$Z_{f}=\{(x,\xi)\in X\times H^{0}(A,\Omega^{1}_{A}) ~|~~ f^{*}\xi(x)=0 \}$$ and $$S_{f}\coloneqq(f\times \mathrm{id})(Z_{f})=\{(a,\xi)\in A\times H^{0}(A,\Omega^{1}_{A}) ~|~~  f^{*}\xi(x)=0,\ \text{for some }x\in f^{-1}(a) \}.$$
Note that the  map $\phi\colon S_f\cap \{(A-C)\times H^0(A, \Omega_A^1)\}\to H^0(A, \Omega_A^1)$ induced by the second projection is dominant, since for a general holomorphic 1-form $\xi$ on $A$, $\xi|_Y$ has isolated zeros on $Y-C$. By contradiction hypothesis, the general fibre of $\phi$ has dimension $\geq 1$, it follows that $$\dim S_f\geq \dim S_f\cap \{(A-C)\times H^0(A, \Omega_A^1)\}\geq \dim H^0(A, \Omega_A^1)+1.$$ However, $\dim S_f\leq \dim A$ according to Lemma~\ref{lem:conormalspace-nonlinear}. This is a contradiction. We are done for Step 4.


\textbf{Step 5. We show that $W^3(X)$ is not linear.}

For any holomorphic 1-form $\omega\in \cS_1$, $\pi^{\ast}(\omega|_Y)\in W^3(X)$ since $\pi^{\ast}(\omega|_Y)$ has zeros $E\times\{y\}$ for some point $y\in C$ such that $\omega|_Y(y)=0$. We have seen that $k^\ast\colon H^0(A, \Omega_A^1)\to H^0(Y, \Omega_Y^1)$ is an isomorphism and $\pi^{\ast}\colon H^0(Y, \Omega_Y^1)\to H^0(X, \Omega_X^1)$ is injective by the Lefschetz hyperplane theorem. Then  $\pi^{\ast}k^\ast\cS_1$ is irreducible, not linear, and $\pi^{\ast}k^\ast\cS_1\subseteq W^3(X)$.  Moreover, since $\cS_1$ is not contained in any proper vector subspace of $H^0(A,\Omega_A^1)$, $\pi^{\ast}k^\ast\cS_1$ is not contained in any proper vector subspace of $\pi^{\ast}k^\ast H^0(A,\Omega_A^1)$.

Suppose by contradiction that $W^3(X)$ is linear. By Equation (\ref{eq:w1-from-Y}) and previous discussion, we have $$\pi^{\ast}k^\ast\cS_1\subseteq W^3(X)\subseteq\pi^{\ast}k^\ast H^0(A,\Omega_A^1)=\pi^{\ast}H^0(Y,\Omega_Y^1).$$  %which is a 5-dimensional vector space. 
Therefore, we get 
\begin{equation}\label{eq:w3=wholespace}
    W^3(X)=\pi^{\ast}k^\ast H^0(A,\Omega_A^1)=\pi^{\ast}H^0(Y,\Omega_Y^1).
\end{equation} 
By Step 4 and Equation (\ref{eq:w3=wholespace}), for a general holomorphic 1-form $\xi$ on $A$, $\pi^{\ast}(\xi|_Y)$ has higher-dimensional zeros $Z(\pi^{\ast}(\xi|_Y))$ in $E\times C$. 

\textbf{Claim $\spadesuit$: $\pi (Z(\pi^{\ast}\xi|_Y))\neq C$ for general $\xi$ on $A$.} Suppose by contradiction that $\pi (Z(\pi^{\ast}\xi|_Y))=C$  for general $\xi$ on $A$. Consider $$S_{f}\coloneqq(f\times \mathrm{id})(Z_{f})=\{(a,\xi)\in A\times H^{0}(A,\Omega^{1}_{A}) ~|~~  f^{*}\xi(x)=0,\ \text{for some }x\in f^{-1}(a) \},$$
and denote  by $\psi\colon S_f\to H^0(A,\Omega_A^1)$ the natural projection. By the contradiction hypothesis, we have the general fibre of $\psi$ has dimension $\geq 1$, it follows that $$\dim S_f \geq \dim H^0(A, \Omega_A^1)+1.$$ However, $\dim S_f\leq \dim A$ according to Lemma~\ref{lem:conormalspace-nonlinear}. This is a contradiction. Claim $\spadesuit$ holds true.

We proceed by analyzing the tangent maps. Taking any closed point $x=(e,y)\in E\times C$, via the commutative diagram (\ref{dia:full-construction}), we have the induced commutative diagram of tangent maps \begin{equation}\label{dia:tangent maps}
\xymatrix{
 T_eE\oplus T_yC \ar[r]^-{p'_{*,x} }\ar@{^{(}->}[d]_{\beta_{*,x}}  \ar@/_4 pc/[dd]_{\alpha_{*,x}} &T_yC\ar@{^{(}->}[d]^{\gamma_{*,y}}\\
   T_eE\oplus T_yY \ar[r]^-{p_{*,x}} & T_yY \\  
	T_xX\ar[ru]_{\pi_{*,x}}    \ar@{_{(}->}[u]^{j_{*,x}},
}
\end{equation} 
where the inclusions $\alpha_{*,x}, \beta_{*,x}, \gamma_{*,y}, j_{*,x}$ are induced by natural inclusions, and the projections $p'_{*,x}, p_{*,x}$ are induced by the second projections. From commutative diagram (\ref{dia:tangent maps}), we know that the tangent map $\pi_{*,x}\colon T_xX\to T_yY$ factors through $T_eE\oplus T_yY$. Since $j_{*,x}$ is injective and $p_{*,x}$ only kills $T_eE$, we obtain that 
\begin{equation}\label{eq:ker-3dim}
    \ker\pi_{*,x}=T_eE,\ \text{and}\ \dim \pi_{*,x}(T_xX)=3.
\end{equation} 
For a general 1-form $\xi$ on $A$, $\pi^{\ast}(\xi|_Y)$ has 1-dimensional zeros of type $E_y=E\times \{y\}$ for some closed points $y\in C$ by Claim $\spadesuit$.  Note that $$\pi^*(\xi|_Y)(E_y)=0\Leftrightarrow \xi|_Y(\bigcup_{x\in E_y}\pi_{*,x}(T_xX))=0.$$ According to Equation (\ref{eq:ker-3dim}), we have the following two situations. 

1) If $\bigcup_{x\in E_y}\pi_{*,x}(T_xX)=T_yY$, we know that $\xi|_Y$ kills $T_yY$ if and only if $\xi|_Y$ has isolated zeros at $y\in C$. This implies that $\xi\in\cS_1$. Denote $$\Sigma_1\coloneqq \{y\in C\ |\ \bigcup_{x\in E_y}\pi_{*,x}(T_xX)= T_yY\}.$$ Then we have $$\dim\{\xi\in H^0(A, \Omega_A^1)\ |\ \pi^{\ast}(\xi|_Y)(E_y)=0\ \text{for some}\ y\in \Sigma_1\}\leq 2.$$

2) If $\bigcup_{x\in E_y}\pi_{*,x}(T_xX)\subsetneq T_yY$. Consider the exact conormal  sequence

\[
    \begin{tikzcd}
        0 \arrow[r]&\cN_{Y/A,y}^{*}\arrow[r]&\{y\}\times H^0(\Omega_A^1)\arrow[r, "\mu"]&  \Omega_{Y,y}^1 \arrow[r]&0
    \end{tikzcd}
    \]
    where $\mu(\xi)=\xi|_Y$.  For any fixed $y_0\in C$ such that $\bigcup_{x\in E_{y_0}}\pi_{*,x}(T_xX)\subsetneq T_{y_0}Y$, the holomorphic 1-forms $\xi$ on $A$ such that $\pi^{\ast}(\xi|_Y)$ have zeros along $E_{y_0}$ form a 2-dimensional cone, since $\dim \bigcup_{x\in E_y}\pi_{*,x}(T_xX)=3$ by Equation (\ref{eq:ker-3dim}). 
    Denote $$\Sigma_2\coloneqq \{y\in C\ |\ \bigcup_{x\in E_y}\pi_{*,x}(T_xX)\subsetneq T_yY\}.$$ Then we have $$\dim\{\xi\in H^0(A, \Omega_A^1)\ |\ \pi^{\ast}(\xi|_Y)(E_y)=0\ \text{for some}\ y\in \Sigma_2\}\leq 3.$$
    
All in all, this contradicts the previous conclusion that for a general holomorphic 1-form $\xi$ on $A$, $\pi^{\ast}(\xi|_Y)$ has zeros along $E_y$ for some $y\in C$. We are done for the proof of this theorem.
\end{proof}














\subsection{Basic properties of non-formal 1-forms}

With the previous non-linear example, one could expect that holomorphic 1-forms that have zeros and do not come from cohomology jump loci are quite wild. We will discuss this kind of  holomorphic 1-forms in this subsection.

\begin{definition}
Let $X$ be a smooth complex projective variety. We call a holomorphic 1-form $\omega\in W^i(X)\backslash V^i(X)$ for some $i$ a \textit{non-formal}\footnote{The theory of cohomology jump loci  of variety $X$ is essentially derived from the formality property of the underlying CW complex of $X$, so we name the holomorphic 1-forms that have zeros and do not come from cohomology jump loci in this way.} holomorphic 1-form of $X$.
\end{definition}
 
We have the following criterion for non-formal holomorphic 1-forms, which is a direct consequence of \cite{GL87}.
\begin{proposition}\label{nonHodge}
Let $\omega$ be a holomorphic 1-form with zeros on $X$ such that $\codim_XZ(\omega)=k$. Then the following statements are equivalent :
    \begin{itemize}
        \item[(a)] $H^{0}(X,\cH^k(\Omega_X^{\bullet},\wedge\omega))=0$, where $\cH^k(\Omega_X^{\bullet},\wedge\omega)$ is the $k$-th cohomology sheaf of the Koszul complex $(\Omega_X^{\bullet},\wedge\omega)$ of sheaves of holomorphic forms;
        \item[(b)] $\omega$ is a non-formal holomorphic $1$-form.
    \end{itemize}
\end{proposition}
\begin{proof}
   Denote the Koszul complex $(\Omega_X^{\bullet},\wedge\omega) $  by $K^{\bullet}(\omega)$, and the $k$-th cohomology sheaf of $K^{\bullet}(\omega)$  by $\cH^k$. Then we have two  spectral sequences:

\[
^{'}E_{1}^{p,q} = H^{q}(X, K^{p}(\omega)) \Rightarrow \HH^{p+q}(K^{\bullet}(\omega))
\]

\[
^{''}E_{2}^{p,q} = H^{p}(X, \cH^{q}) \Rightarrow \HH^{p+q}(K^{\bullet}(\omega))
\]
Note that we have following facts \cite[Proposition~3.7 and Lemma~3.8]{GL87}:
\begin{itemize}
    \item $\Supp \cH^{i}\subset Z(\omega)$ for all $i$;
    \item $^{'}E_{1}^{p,q}$ degenerates at  $E_{2}$ page;
    \item $\cH^{i}= 0$ for $i < k$.
\end{itemize} 
Then based on the spectral sequence $^{''}E_{2}^{p,q}$, we have 
\[
^{''}E_{2}^{0,k}= H^{0}(X, \cH^{k})=0\text{ if and only if }   \HH^k(K^{\bullet}(\omega))=0.
\]
From the spectral sequence of $^{'}E_{1}^{p,q}$, we infer that
 $\HH^k(K^{\bullet}(\omega))=0$ if and only if the sequence
 \[
 \cdots\rightarrow H^i(X,\Omega_{X}^{k-i-1}) \xrightarrow{\wedge \omega} H^i(X,\Omega_{X}^{k-i}) \xrightarrow{\wedge \omega} H^i(X,\Omega_{X}^{k-i+1}) \xrightarrow{\wedge \omega} \cdots 
 \]
 is exact at $H^i(X,\Omega_{X}^{k-i})$ for all $i$, via Hodge decomposition, which is also equivalent to that the following sequence
 \[
 \cdots\rightarrow H^{k-1}(X,\CC) \xrightarrow{\wedge \omega} H^{k}(X,\CC) \xrightarrow{\wedge \omega} H^{k+1}(X,\CC)\xrightarrow{\wedge \omega} \cdots 
 \]
 is exact at $H^k(X,\CC)$. Then the equivalence holds true.
\end{proof}

Based on the above cohomological criterion, we obtain the following simple lemma.
\begin{lemma}\label{lem:positivedim}
    If $\omega$ is a non-formal 1-form, then $\dim Z(\omega)>0$.
\end{lemma}
\begin{proof}
     Suppose that $\dim X=n$. In local coordinates $x_1,x_2,\cdots,x_n$, the holomorphic 1-form $\omega$ can be written into $\omega=\sum_{i=1}^{n}f_idx_i$ locally. Consider the map by wedging $\omega$ :
$$
\begin{tikzcd}
\Omega_{X}^{n-1} \arrow[r,"\wedge \omega"] & \Omega_{X}^{n} \\
\eta \arrow[r]                 & \sum_{i=1}^{n}f_idx_i\wedge\eta
\end{tikzcd}
$$
By \cite[Lemma~4.2]{Sch21}, $\cH^n(\Omega_X^{\bullet},\wedge\omega)\cong \Omega_X^n|_{Z(\omega)}$. Suppose by contradiction that $\dim Z(\omega)=0$. By Proposition~\ref{nonHodge}, we would have 
$$
H^0(X, \cH^n)=H^0(Z(\omega), \Omega_X^n|_{Z(\omega)})=0,
$$
which is a contradiction.
\end{proof}
 

In the rest of this section, we provide some discussions on surfaces with non-formal 1-forms. 


\begin{lemma}\label{lem:nonreduced}
    Let $S$ be a smooth projective surface, $f\colon S\to E$ be a morphism onto an elliptic curve $E$ and $\omega\in H^0(E, \Omega_E^1)$ be a nonzero holomorphic 1-form. If $f^{\ast}\omega$ is a non-formal holomorphic 1-form, then there exists a 1-dimensional irreducible component $F$ of $Z(\omega)$, such that $F$ is $f$-vertical and the component of the scheme-theoretic fibre $f^{\ast}(f(F))$ supported on $F$ is non-reduced. 
\end{lemma}
\begin{proof}
    By Lemma~\ref{lem:positivedim}, there exists $F\subset Z(\omega)$ with $\dim F=1$. Since $f$ is generically smooth, $F$ is $f$-vertical, say $f(F)=p$ for some closed point $p\in E$. Now consider the local homomorphism between discrete valuation rings $f^{\sharp}\colon \cO_{E,p}\to \cO_{S,\eta}$ where $\eta$ is the generic point of $F$. Fix  uniformizers $t\in \cO_{E,p}\text{ and } s\in \cO_{S,\eta}$ respectively, and denote $f^{\sharp}(t)=s^{k}$ for some integer $k$. For any closed $x\in F$, the cotangent map $\mathfrak{m}_{p}/{\mathfrak{m}_{p}^2}\to\mathfrak{m}_{x}/{\mathfrak{m}_{x}^2}$ is zero, where $\mathfrak{m}$ denote maximal ideals. Then $\mathfrak{m}_{p}/{\mathfrak{m}_{p}^2}\to\mathfrak{m}_{\eta}/{\mathfrak{m}_{\eta}^2}$ is the zero map. Therefore, $f^{\sharp}(\mathfrak{m}_p)\subset {\mathfrak{m}_{\eta}^2}$, that is $k\geq 2$.
\end{proof}
\begin{remark}\label{rmk:q=1}
    If $\alb.\dim S=1$, then the converse is also true. Indeed, we claim that $q(S)=1$. Granting this claim for the moment. The Albanese morphism $a_S$ of $S$ is a fibration over an elliptic curve $\Alb(S)$. Then for $0\neq\omega\in H^0(\Alb(S), \Omega^1_{\Alb(S)})$ the sequence  \[
 0\to H^0(S,\CC) \xrightarrow{\wedge a_S^{\ast}\omega} H^1(S,\CC) \xrightarrow{\wedge a_S^{\ast}\omega} H^2(S,\CC) 
 \]
        is exact at $H^1(S,\CC)$, since $q(S)=1$.  Also, $\dim Z(f^{\ast}\omega)=1$. Hence $f^{\ast}\omega$ is non-formal by definition.
        
        As for the claim. If $q(S)\geq 2$. Then the Albanese fibration $a_S\colon S\to C$ is a fibration over a smooth curve $C$ of genus $q(S)$. Then there is a holomorphic 1-form $\eta$ on $C$ which is not proportional to $\omega$, such that $\eta\wedge\omega=0$ on $C$. In this case, the sequence  \[
 0\to H^0(S,\CC) \xrightarrow{\wedge a_S^{\ast}\omega} H^1(S,\CC) \xrightarrow{\wedge a_S^{\ast}\omega} H^2(S,\CC) 
 \]
        is not exact at $H^1(S,\CC)$, which is a contradiction.
        
\end{remark}
   

\begin{lemma}\label{lem:exist-nonhodge+albdim1}
    Let $X$ be a smooth projective variety admitting a non-formal 1-form $\omega$ and $\alb.\dim(X)=1$, then the irregularity $q(X)=1$.
\end{lemma}
\begin{proof}
    Let $\alpha\colon X\to \Alb(X)$ be the Albanese morphism. By assumption, we know that $\alpha(X)$ is a smooth curve $C$ of genus $q(X)$ (see e.g., \cite[Proposition 9.19]{Ueno75}).  Suppose by contradiction that $q(X)=g(C)>1$, then there exists another holomorphic 1-form $\omega'$ linearly independent with $\omega$ in $H^0(X,\Omega_X^1)$ such that $\omega\wedge \omega'=0$. This implies the complex $$0\to H^0(X, \mathcal{O}_X)\overset{\wedge\omega}{\longrightarrow} H^0(X,\Omega^1_X)\overset{\wedge\omega}{\longrightarrow} H^0(X,\Omega^2_X)$$ is not exact at $H^0(X,\Omega^1_X)$. However, since $\omega$ is non-formal with $d=\textnormal{codim} Z(\omega)\geq 1$, $$0\to H^0(X, \mathcal{O}_X)\overset{\wedge\omega}{\longrightarrow} H^0(X,\Omega^1_X)\overset{\wedge\omega}{\longrightarrow} H^0(X,\Omega^2_X)$$  is exact at $H^0(X,\Omega^1_X)$. This is a contradiction.
\end{proof}

\begin{proposition}
Let $S$ be a smooth projective surface with $\alb.\dim(S)=1$. Then $S$ admits a non-formal 1-form, if and only if $q(S)=1$ and there exists a fibre of the Albanese map of $S$ admitting a non-reduced component.
\end{proposition}

\begin{proof}
    This follows directly from Lemma \ref{lem:nonreduced}, Remark \ref{rmk:q=1}, and Lemma \ref{lem:exist-nonhodge+albdim1}.
\end{proof}



%\begin{corollary}
 %   Let $S$ be a smooth projective surface with $q(S)=1$, then the Albanese morphism $a_S$ either is smooth or has a singular fibre with a non-reduced component.
%\end{corollary}
%\begin{proof}
 %   If $a_S$ is not smooth, then $Z(f^{\ast}\omega)\neq\emptyset$. Also, the sequence \[
 %0\to H^0(S,\CC) \xrightarrow{\wedge a_S^{\ast}\omega} H^1(S,\CC) \xrightarrow{\wedge a_S^{\ast}\omega} H^2(S,\CC) 
 %\]
  %        is exact at $H^1(S,\CC)$. Hence $f^{\ast}\omega$ is non-Hodge. By Lemma~\ref{lem:nonreduced}, there is a non-reduced 1-dimensional component $F$ contained in $Z(f^{\ast}\omega)$.
%\end{proof}

%{\color{red}This corollary should hold for higher varieties.}


\begin{proposition}
    Let $S$ be a smooth minimal projective surface admitting a non-formal 1-form $\omega$ and $\alb.\dim(S)=2$, then $S$ must be of general type.
\end{proposition}
\begin{proof}
    Since $S$ is of maximal Albanese dimension, then the Kodaira dimension $\kappa(S)\geq 0$ and $q(S)\geq 2$. Suppose by contradiction that $S$ is not of general type, then $K_S^2=0$. The assumption that $\omega$ is non-formal implies that the sequence \[
 0\to H^0(S,\CC) \xrightarrow{\wedge \omega} H^1(S,\CC) \xrightarrow{\wedge \omega} H^2(S,\CC) 
 \]
        is exact at $H^1(S,\CC)$. It follows that there exists $\eta\in H^0(S, \Omega_S^1)$, not proportional with $\omega$, such that $\eta\wedge\omega\neq 0$ in $H^0(S, \Omega_S^2)$. Also, there exists a 1-dimensional component $F\subset Z(\omega)$. Then we have  $F\subset Z(\eta\wedge\omega)$. We infer that $aF+M$ is linearly equivalent to $K_S$ for some effective divisor $M$ on $S$ and positive integer $a$. In the case, $0=K_S^2=aK_SF+K_SM$ and $K_S$ is nef, thus $K_SF=0$. 
        
        We claim that $F^2<0$. Indeed, if $F^2\geq 0$, by \cite[Theorem~1]{Spu88}, $F$ is a rational multiple of a fibre of an irrational pencil $f\colon S\to C$ and $\omega\in f^{\ast}H^0(C, \Omega_C^1)$. Hence it must be the whole fibre. But the arithmetic genus of $F=1$, which implies general fibre of $f$ is elliptic curve. Since $\omega$ is non-formal, then $g(C)=1$. In this case, $2\leq q(S)\leq g(C)+1=2$, which implies that $S$ is a  product of elliptic curve by the Lemme in the Appendix of \cite{Deb82}, that is a contradiction since such a surface doesn't contain any non-formal 1-form. 
        
        Hence by genus formula, $F$ is a $(-2)$-curve. Since $K_S^2=0$, then the Stein factorization of the canonical morphism of $S$ yields an elliptic fibration $\phi\colon S\to B$ with a general fibre $E$. $K_SF=0$ implies that $F$ is contained in a singular fibre of $\phi$. Since $\phi$ has a singular fibre with a $(-2)$-component $F$,  $S$ is not isomorphic to $E\times B$. Hence $q(S)=g(B)$ by Lemme in the Appendix of \cite{Deb82}. By the argument in Remark~\ref{rmk:q=1}, the base curve $B$ must be elliptic and hence $q(S)=g(B)=1$, which is impossible.
\end{proof} 
\begin{proposition}
    Let $S$ be a surface with a non-formal 1-form, then $S$ is not dominated by a product of curves of higher genera. In particular, $S$ is not a product-quotient surface.
\end{proposition}
    
 
\begin{proof}
    Suppose that $f\colon C_1\times C_2\to S$ is a surjective morphism, where $C_1,C_2$ are smooth curves of genera $g(C_1)\geq 2$ and $g(C_2)\geq 2$, respectively. If $\omega\in H^0(S, \Omega_S^1)$ is a non-formal holomorphic 1-form, then there is a 1-dimensional component  $D\subset Z(\omega)$ and $D^2<0$. It follows that $\dim Z(f^{\ast}\omega)=1$ and $f^{\ast}D\subset Z(f^{\ast}\omega)$. Hence there is an irreducible component $\widetilde{D}\subset f^{\ast}D$ with $\widetilde{D}^2<0$, which is a contradiction since the higher dimensional zero locus of a holomorphic 1-form on $C_1\times C_2$ is either horizontal or vertical with respect to  projections to its factors, which has non-negative self-intersection numbers.
\end{proof}

































 

%references added in the ref.bib file
\bibliographystyle{plain}
\bibliography{ref}

% references added in the bibitem form
%\begin{thebibliography}{99}
%\bibitem[BL]{BL}Birkenhake, Christina; Lange, Herbert.  Complex abelian varieties. Grundlehren Math. Wiss., 302[Fundamental Principles of Mathematical Sciences]Springer-Verlag, Berlin, 2004.
%\end{thebibliography}



\end{document}
